\section{Complete CM Grounding Proof}\label{app:structure}

This appendix supplies the detailed statements underlying
Theorem~\ref{main:cm}. We use $A_S^\beta$ from
\eqref{eq:mainmatrices}, the admissible seed lattice from
Section~\ref{sec:prelim}, and $a=\mu(\{0\})=f(\infty)$.
The diminishing-drop inequality is
\begin{equation}\label{eq:super}
\Phi(S)-\Phi(S+v)\geq\Phi(T)-\Phi(T+v),\quad S\subseteq T,\ v\notin T.
\end{equation}
\begin{lemma}[One coupling for all grounding sets]\label{lem:coupling}
There is a coupling, independent of the selected set, with killing times
$\mathsf T_S$ such that $\mathsf T_{S\cup T}=\min(\mathsf T_S,\mathsf T_T)$.
For soft grounding,
\[
 (e^{-tA_S^\beta})_{ij}=\Prob_i(\mathsf X_t=j,\ t<\mathsf T_S),\quad t\geq0.
\]
For hard grounding the same identity holds after embedding
$e^{-tA_S^\infty}$ into an $n\times n$ matrix by zeros in rows and
columns in $S$, including at $t=0$.
\end{lemma}
\begin{proof}
Let $\mathsf X$ be the continuous-time walk on $V$ with jump rate $w_{ij}$ from
$i$ to $j$, holding rate $d_i=\sum_j w_{ij}$, and generator $-L$.
It exists for all time since the state space and all rates are finite;
an isolated state has infinite holding time. Independently sample
$E_0\sim\operatorname{Exp}(\sigma)$ (take $E_0=\infty$ if $\sigma=0$),
and $E_v\sim\operatorname{Exp}(\beta)$ independently for every $v$.
Define occupation times $\ell_v(t)=\int_0^t\ind_{\{\mathsf X_u=v\}}du$ and
\[
 \kappa_v=\inf\{t:\ell_v(t)\geq E_v\},\qquad
 \mathsf T_S=E_0\wedge\min_{v\in S}\kappa_v.
\]
These clocks run in local occupation time. In particular an unvisited
vertex cannot kill the walk. The displayed minimum identity holds
pathwise, simultaneously for all sets. Conditioning on the walk gives
\[
 \Prob(t<\mathsf T_S\mid \mathsf X)=
 e^{-\sigma t}\exp\left(-\beta\sum_{v\in S}\ell_v(t)\right).
\]
The exponential memoryless property implies that while at $i$ the
additional killing rate is $\sigma+\beta\ind_S(i)$. The transient
transition matrix $K_S(t)$ therefore has initial value $I$ and generator
$Q_{ij}=w_{ij}$ for $i\ne j$, $Q_{ii}=-d_i-\sigma-\beta\ind_S(i)$.
The Markov property and the first-step expansion give $K'_S(t)=K_S(t)Q$;
uniqueness of this finite linear ODE yields $K_S(t)=e^{-tA_S^\beta}$.

For hard grounding replace $\kappa_v$ by the hitting time
$H_v=\inf\{t\geq0:\mathsf X_t=v\}$ and keep the same $E_0$ and base walk.
On $U=V\setminus S$ the transient rates are $w_{ij}$, $i,j\in U$, $i\ne j$,
and diagonal $-d_i-\sigma$; jumps into $S$ cause absorption.
Thus the transient generator is $-M[U,U]$. A start in $S$ has
$\mathsf T_S=0$ and contributes zero even at $t=0$. A surviving walk cannot
end in $S$. These facts give precisely the zero embedding and prove
the identity for every $i,j$. All the constructions use one base walk
and the same independent clocks for every set.
\end{proof}

\begin{theorem}[CM grounding traces]\label{thm:cm}
For every CM $f$ and $\beta\in(0,\infty]$, $\Phi_f^\beta$ is
nonincreasing and supermodular on the admissible lattice.
For each admissible $S$ and $v\notin S$, the drop is strictly positive
in the soft case exactly when $f$ is nonconstant, and in the hard case
exactly when $f\not\equiv0$.
The soft-to-hard limit is
\begin{equation}\label{eq:hardlimit}
 \lim_{\beta\to\infty}\Phi_f^\beta(S)
 =\Phi_f^\infty(S)+|S|a.
\end{equation}
In particular, the two objectives agree in the limit when $f(\infty)=0$.
\end{theorem}
\begin{proof}
Fix a starting vertex $i$, time $t$, and one coupled outcome. Define
\[
 Z_S=\ind_{\{\mathsf X_t=i\}}\ind_{\{t<\mathsf T_S\}}.
\]
For $S\subseteq T$, $Z_S\geq Z_T$. For the added vertex clock
$\kappa_v$ (or hitting time in the hard case),
\[
 Z_S-Z_{S\cup\{v\}}
 =\ind_{\{\mathsf X_t=i\}}\ind_{\{t<\mathsf T_S\}}\ind_{\{\kappa_v\leq t\}}.
\]
Replacing $S$ by $T$ only decreases this product. Taking expectations
and summing over all $i\in V$ proves monotonicity and
\eqref{eq:super} for $h_S(t)=\tr e^{-tA_S^\beta}$, using the zero
embedding in the hard case. This argument includes $t=0$: then the
soft heat trace is $n$ and the hard heat trace is $n-|S|$.

Diagonalizing each positive definite matrix and using Tonelli's theorem
for the nonnegative scalar integrands gives
\[
 \Phi_f^\beta(S)=\int_{[0,\infty)}h_S(t)\,\mu(dt)
 =\sum_i\int_{[0,\infty)}\E_i Z_S\,\mu(dt).
\]
The measure is sigma-finite, since $\mu([0,j])\leq e^j f(1)$.
Every integral is finite by the finite spectral sum. Differences can
therefore be taken, or equivalently Tonelli can be applied directly
to the nonnegative drops. The pointwise drop inequalities integrate
to the claimed inequalities without subtracting infinities.

For strictness, start at $v\notin S$. The event of no jump before $t$
and survival of base killing has probability $e^{-(d_v+\sigma)t}$.
For soft grounding it contributes
$e^{-(d_v+\sigma)t}(1-e^{-\beta t})>0$ to the heat-trace drop for
every $t>0$. Thus any mass of $\mu$ on $(0,\infty)$ gives a strict drop;
if there is no such mass $f=a$ and the soft objective is $na$.
The representation shows that $f$ is constant exactly in this latter
case. For hard grounding the no-jump event contributes
$e^{-(d_v+\sigma)t}>0$, and the drop at $t=0$ is one. Hence every
nonzero $\mu$ gives a strict drop; $\mu=0$ gives the zero objective.

For the limit use common unit exponential variables and set
$E_v^{\beta}=E_v^1/\beta$. For fixed $t>0$, the conditional survival
factor converges to the indicator that no selected vertex has positive
occupation time before $t$. Apart from a null event of a jump exactly
at $t$, this is hard survival; a start in $S$ also has positive
occupation time. Bounded convergence gives
$h_S^\beta(t)\to h_S^\infty(t)$ for $t>0$.
At $t=0$ the limit remains $n$, whereas the hard trace is $n-|S|$.
Choose a fixed $\beta_0>0$. For $\beta\geq\beta_0$,
$0\leq h_S^\beta(t)\leq h_S^{\beta_0}(t)$, an integrable bound by
admissibility. Dominated convergence proves \eqref{eq:hardlimit},
including $S=V$.
\end{proof}

\begin{corollary}[Powers, heat, and current-flow closeness]\label{cor:instances}
The theorem applies to $f(x)=x^{-s}$ for every $s>0$, to
$f(x)=e^{-\tau x}$ for $\tau\geq0$, and to $(x+c)^{-1}$ for $c\geq0$.
For a connected graph, $\sigma=0$, and $\varnothing\ne S\subsetneq V$,
the group current-flow closeness is
$\operatorname{CFCC}(S)=n/\tr L_{-S}^{-1}$.
Thus the reciprocal-CFCC supermodularity of Li et al.~\cite{li2019}
is the hard-grounding $s=1$ case.
\end{corollary}
\begin{proof}
The positive representations are, respectively,
\begin{align*}
 x^{-s}&=\frac1{\Gamma(s)}\int_0^\infty t^{s-1}e^{-tx}dt,\\
 e^{-\tau x}&=\int e^{-tx}\delta_\tau(dt),\\
 (x+c)^{-1}&=\int_0^\infty e^{-tx}e^{-ct}dt.
\end{align*}
The first follows by $u=tx$ in the defining Gamma integral, and the
third is an elementary exponential integral. Apply Theorem~\ref{thm:cm}.
For the electrical identity, ground $S$ and inject a unit current at
$u\notin S$. Kirchhoff's equations on $U=V\setminus S$ are
$L[U,U]p=e_u$. Thus the source voltage, equal to its effective
resistance to $S$, is $(L[U,U]^{-1})_{uu}$. Summing these voltages
and using the definition of group current-flow closeness gives the
formula. Multiplication by $1/n$ preserves supermodularity of the
reciprocal objective.
\end{proof}

\begin{proposition}[Edge addition has a different structure]\label{prop:edge}
Even for $f(x)=x^{-1}$, edge addition can violate supermodularity on a
connected five-vertex graph.
\end{proposition}
\begin{proof}
Take the unit-weight path $2-1-0-4-3$, $\sigma=1/10$, and candidate
unit edges $e=\{0,2\}$, $g=\{0,3\}$. Write
$H(B)=\tr(L+\sum_{b\in B}b_b b_b^\top+I/10)^{-1}$, where
$b_{\{u,v\}}=e_u-e_v$. Exact rational inversion gives
\begin{align*}
 H(\varnothing)&=\frac{966250}{72181},&
 H(\{e,g\})&=\frac{204350}{17391},\\
 H(\{e\})&=H(\{g\})=\frac{2486050}{197501}.&&
\end{align*}
Rational Gaussian elimination gives these values; the inverse
identities are verified by direct matrix multiplication.
Their mixed difference is exactly
\[
 H(\{e,g\})-H(\{e\})-H(\{g\})+H(\varnothing)
 =-\frac{9866618000}{257984349711}<0.
\]
For grounding the off-diagonal transition rates remain fixed and
only killing changes. Adding an edge instead changes two transition
rates, hence changes the base path and its return event. Additional
jumps cannot be represented by deleting trajectories with independent
vertex killing clocks. Algebraically the edge direction $e_u-e_v$
has both signs, unlike the grounding direction $e_v$; for $R=A^{-1}$,
direct differentiation gives
$\partial_x\partial_y\tr(A+xbb^\top+ycc^\top)^{-1}|_{0}
=2(b^\top Rc)(b^\top R^2c)$, which has no fixed sign.
Indeed $\partial_x\tr R=-b^\top R^2b$ and
$\partial_yR=-Rcc^\top R$ give the displayed derivative by the product
rule. This identifies exactly the unavailable step in the coupling
proof for edge addition.
\end{proof}



\section{Hardness Proofs}\label{app:hardness}

The first four results prove the inverse-power and heat parts of
the selection-hardness claims discussed in Section~\ref{sec:structure},
together with the finite-penalty variants.
Here the reduction's $k$ denotes its specified total cardinality;
a singleton-seed call uses $k-1$ additional choices.
The second part proves the stronger exact-budget spectral barrier.
\begin{lemma}[Cubic vertex-cover gap]\label{lem:cover}
Let $G$ be a connected unit-weight $3$-regular graph, let
$\varnothing\ne S\subsetneq V$, put $r=n-|S|$, and let $e_U$ count the
edges induced by $U=V\setminus S$. For every $s>0$ and $\sigma\geq0$,
\begin{equation}\label{eq:covergap}
\begin{aligned}
 \tr (L_{-S}+\sigma I)^{-s}
 &\geq r(3+\sigma)^{-s}+\kappa_{s,\sigma}^{\rm curv}e_U,\\
 \kappa_{s,\sigma}^{\rm curv}&=\frac{s(s+1)}{(6+\sigma)^{s+2}}>0.
\end{aligned}
\end{equation}
Equality with the baseline $r(3+\sigma)^{-s}$ holds exactly when
$S$ is a vertex cover. Also $\lambda_{\min}(L_{-S}+\sigma I)\geq n^{-2}$.
\end{lemma}
\begin{proof}
Write $C=(3+\sigma)I-\mathsf W_U$, where $\mathsf W_U$ is the induced adjacency
matrix. The eigenvalues of $C$ lie in $(0,6+\sigma]$: positive
definiteness follows from the nonempty Dirichlet boundary, and
$\|\mathsf W_U\|_2\leq3$. For $f(x)=x^{-s}$, throughout this interval
$f''(x)\geq\kappa_{s,\sigma}^{\rm curv}$. Taylor's formula with integral remainder
therefore gives
\[
 f(x)\geq f(b)+f'(b)(x-b)+\tfrac12\kappa_{s,\sigma}^{\rm curv}(x-b)^2,
 \qquad b=3+\sigma.
\]
Sum this inequality over the eigenvalues. The linear terms vanish
because $\tr C=rb$, and the sum of squared deviations is
$\tr \mathsf W_U^2=2e_U$, proving \eqref{eq:covergap}. If $e_U=0$, $C=bI$
gives equality; if $e_U\geq1$ the excess is at least
$\kappa_{s,\sigma}^{\rm curv}$. The equivalence with a vertex cover follows from
the definition.

For the lower eigenvalue bound, extend $x\in\R^U$ by zeros on $S$.
For each $u\in U$ choose a simple path from $u$ to $S$, of length at
most $n-1$. Cauchy--Schwarz bounds $|x_u|^2$ by $n$ times the sum of
edge differences squared along that path, hence by $n x^\top Lx$.
Summing over at most $n$ vertices yields
$\|x\|^2\leq n^2x^\top Lx\leq n^2x^\top(L+\sigma I)x$.
The Rayleigh quotient proves the asserted eigenvalue bound.
\end{proof}

\begin{lemma}[Uniform finite-penalty approximation]\label{lem:softlimit}
Under the hypotheses of Lemma~\ref{lem:cover}, for $0<s<1$ and $B>0$,
\begin{equation}\label{eq:softbound}
\begin{aligned}
 0&\leq \tr(L+\sigma I+B\diag(\ind_S))^{-s}
       -\tr(L_{-S}+\sigma I)^{-s}\\
 &\leq n\left(\frac{1+9n^4}{B}\right)^s.
\end{aligned}
\end{equation}
\end{lemma}
\begin{proof}
Order the vertices as $U,S$ and write
\begin{gather*}
 L+\sigma I=\begin{pmatrix}C&D\\D^\top&E\end{pmatrix},\qquad
 W=\begin{pmatrix}-C^{-1}D\\I\end{pmatrix},\\
 F=E+BI-D^\top C^{-1}D.
\end{gather*}
Because $L+\sigma I\succeq0$ and $C\succ0$, its Schur complement is
positive semidefinite. For completeness, its quadratic form at $y$
is the quadratic form of $L+\sigma I$ at $(-C^{-1}Dy,y)$.
Consequently $F\succeq BI$. Block multiplication verifies
\[
 (L+\sigma I+B\diag(\ind_S))^{-1}
 =J+WF^{-1}W^\top,\qquad J=\begin{pmatrix}C^{-1}&0\\0&0\end{pmatrix}.
\]
Thus $0\preceq A^{-1}-J\preceq\epsilon I$, where
\[
 \epsilon=\frac{1+\|C^{-1}D\|_2^2}{B}
 \leq\frac{1+9n^4}{B}.
\]
Here $\|C^{-1}\|_2\leq n^2$ by Lemma~\ref{lem:cover}, and
$\|D\|_2\leq3$ since it is a submatrix of the cubic adjacency matrix.
The ordered eigenvalues satisfy
$\lambda_j(J)\leq\lambda_j(A^{-1})\leq\lambda_j(J)+\epsilon$
by the min--max principle. For $x\geq0$ and $0<s<1$,
$0\leq(x+\epsilon)^s-x^s\leq\epsilon^s$:
the difference is decreasing in $x>0$, since $s-1<0$, and its
maximum is attained at $x=0$. Apply this scalar inequality to the
matched eigenvalues and sum all $n$ terms, including the $|S|$ zeros
of $J$. Since $\tr J^s=\tr C^{-s}$, this proves \eqref{eq:softbound}.
\end{proof}

\begin{theorem}[Hardness for fixed fractional exponent]\label{thm:hardness}
For each fixed $s\in(0,1)$, finding an exact minimizing set of total
cardinality $k$ is NP-hard in each of the following regimes:
\begin{enumerate}
\item hard grounding on connected unit-weight cubic graphs, with any
fixed $\sigma\geq0$;
\item soft grounding on the same graphs, with any fixed $\sigma\geq0$
and an input integer penalty $\beta$ of polynomial magnitude;
\item soft grounding with $\sigma=0$ and any fixed positive rational
penalty $\beta_0$, on connected rationally weighted cubic graphs.
\end{enumerate}
For $\sigma=0$, the specified-singleton-seed formulation is also
NP-hard under polynomial-time Turing reductions.
\end{theorem}
\begin{proof}
Use vertex cover on connected $3$-regular graphs, the NP-complete
source problem used in Section 4 of Li et al.~\cite{li2019}.
The nontrivial instances have $1\leq k<n$. A cover of size at most
$k$ exists exactly when one of size $k$ exists, since a cover can
be enlarged. Lemma~\ref{lem:cover} proves that whenever a size-$k$
cover exists, every hard-objective minimizer is a cover. An optimal-set
oracle therefore decides vertex cover by returning a set whose edge
coverage can be checked in polynomial time. This proves item 1
without needing an exact representation of a real objective threshold.

For item 2, write $\delta=\kappa_{s,\sigma}^{\rm curv}$. Fix integers
$p>1/s$ and $C>0$ such that $C^{-s}\leq\delta/8$, and choose
\[
 B=C(1+9n^4)n^p.
\]
These integers depend only on the fixed parameters $s,\sigma$.
The chosen $B$ has polynomial magnitude and $O(\log n)$ bits.
By Lemma~\ref{lem:softlimit} the error for every nonempty proper
set is at most
$nC^{-s}n^{-ps}\leq\delta/8$. Put $b=(n-k)(3+\sigma)^{-s}$.
A cover has soft value in $[b,b+\delta/8]$, whereas a noncover
has soft value at least $b+\delta$. Every soft minimizer is therefore
a cover whenever one exists. Again check the returned set's edge
coverage rather than comparing real-valued oracle outputs.

For item 3 use the same construction at $\sigma=0$, and replace
every unit edge weight by $c=\beta_0/B$. The new Laplacian is $cL$,
and exactly
\begin{gather*}
 cL+\beta_0\diag(\ind_S)=c(L+B\diag(\ind_S)),\\
 \tr(cL+\beta_0\diag(\ind_S))^{-s}
 =c^{-s}\tr(L+B\diag(\ind_S))^{-s}.
\end{gather*}
The positive common factor preserves all minimizers. The rational
weights have $O(\log n)$ bits, and the support remains connected and
cubic. This is a polynomial reduction with fixed finite $\beta_0$.

Finally suppose the oracle requires a specified nonempty seed at
$\sigma=0$. Make one call for each singleton seed $\{v\}$ with the
same total cardinality $k$ (or $k-1$ additional choices).
If a size-$k$ cover exists, one such seed belongs to it, and the
separation arguments above ensure that the corresponding oracle
returns a cover. If no cover exists, no returned set passes the edge
coverage test. These $n$ calls are a polynomial-time Turing reduction.
\end{proof}

\begin{corollary}[Strictly convex hard-grounding objectives]\label{cor:hardextended}
For a fixed $C^2$ function $f$ with $f''(x)\geq c_f>0$ on
$(0,6+\sigma]$, the cubic-graph objective obeys
\[
 \tr f(L_{-S}+\sigma I)\geq(n-|S|)f(3+\sigma)+c_f e_U.
\]
Equality with the baseline characterizes vertex covers. Consequently,
hard-grounding optimal-set selection is NP-hard for every fixed $s>0$
with $f(x)=x^{-s}$, and for every fixed $\tau>0$ with
$f(x)=e^{-\tau x}$, at any fixed $\sigma\geq0$.
The singleton-seed qualification in Theorem~\ref{thm:hardness} applies.
\end{corollary}
\begin{proof}
Repeat the Taylor inequality in Lemma~\ref{lem:cover} with $c_f$:
the linear terms sum to zero and the quadratic terms sum to
$c_f\tr \mathsf W_U^2/2=c_f e_U$. A cover has $\mathsf W_U=0$ and achieves the
baseline, whereas a noncover exceeds it by at least $c_f$.
Thus an optimal-set oracle decides the connected cubic vertex-cover
problem exactly as in Theorem~\ref{thm:hardness}. Powers use
$c_f=s(s+1)(6+\sigma)^{-s-2}$ for every $s>0$; heat uses
$c_f=\tau^2e^{-\tau(6+\sigma)}>0$. Checking returned sets avoids
real threshold encoding. The same singleton-seed calls prove the
seeded assertion. At $\tau=0$ the fixed-cardinality heat objective
is constant, so the heat statement requires $\tau>0$.
\end{proof}

The following construction leaves the anchor budget unchanged and
allows every new internal vertex to be selected. For an integer
$\ell\geq2$, define
\begin{equation}\label{eq:pathconstants}
 a_\ell=2-2\cos\frac{\pi}{2\ell-1},\qquad
 b_\ell=2-2\cos\theta_\ell,
\end{equation}
where $\theta_\ell$ is the unique root in
$(0,\pi/(2\ell-1))$ of
\begin{equation}\label{eq:pathroot}
 3\cos((\ell-\tfrac12)\theta)=\cos(\theta/2).
\end{equation}

\begin{lemma}[Parallel-path YES/NO bounds]\label{lem:pathamplification}
Let $H$ be a connected simple cubic graph. Form $G_{m_{\rm path},\ell}$ by
replacing every edge of $H$ by $m_{\rm path}$ internally vertex-disjoint
paths of length $\ell$, with unit edge weights. The original
vertices remain as the paths' endpoints. Then:
\begin{enumerate}
\item If $S$ is an original vertex cover with $|S|<|V(H)|$, then
\begin{equation}\label{eq:pathyes}
 \lambda_{\min}(L(G_{m_{\rm path},\ell})_{-S})
 \geq\frac{a_\ell}{1+2(1+a_\ell)/(3m_{\rm path})}.
\end{equation}
\item Let $T$ be any set of $k$ vertices of $G_{m_{\rm path},\ell}$, including
arbitrary internal vertices. If $T\cap V(H)$ is not a cover and
$m_{\rm path}(\ell-1)\geq8k$, then
\begin{equation}\label{eq:pathno}
 \lambda_{\min}(L(G_{m_{\rm path},\ell})_{-T})
 \leq b_\ell+\frac{24k}{m_{\rm path}(\ell-1)}.
\end{equation}
\item The constants satisfy $0<b_\ell<a_\ell\leq1$,
$b_\ell/a_\ell>c_*$ for every finite $\ell$, and
\begin{equation}\label{eq:pathlimit}
 \lim_{\ell\to\infty}\frac{b_\ell}{a_\ell}
 =c_*:=\left(\frac{2\arccos(1/3)}{\pi}\right)^2.
\end{equation}
\end{enumerate}
\end{lemma}
\begin{proof}
Put $h=\ell-1$. The quadratic form
\[
 q(y)=\sum_{j=1}^{h-1}(y_j-y_{j+1})^2+y_h^2
\]
has smallest eigenvalue $a_\ell$. Indeed its matrix is the path
tridiagonal with first diagonal entry 1, all other diagonal entries
2, and off-diagonals $-1$ (for $h=1$ the matrix is $[1]$).
Its eigenvectors have entries $\cos((j-\tfrac12)\vartheta)$,
with $\vartheta=(2r-1)\pi/(2h+1)$, $r=1,\ldots,h$;
substitution gives the distinct eigenvalues $2-2\cos\vartheta$.
Thus $q(y)\geq a_\ell\|y\|^2$.

When $S$ is a cover, each component after grounding is either a
star with $3m_{\rm path}$ length-$\ell$ rays ending at ground, or an internal
path with both endpoints grounded. For a star let $x$ be its center
value, $y^{(r)}$ the internal values of ray $r$, and write
$Y=\sum_r\|y^{(r)}\|^2$. Its energy is
\[
 E=\sum_{r=1}^{3m_{\rm path}}\bigl((x-y^{(r)}_1)^2+q(y^{(r)})\bigr)
 \geq a_\ell Y.
\]
Averaging $x^2\leq2(x-y^{(r)}_1)^2+2(y^{(r)}_1)^2$ over the
rays yields $x^2\leq2(E+Y)/(3m_{\rm path})$. Hence
\[
 x^2+Y\leq\left(\frac{1+2/(3m_{\rm path})}{a_\ell}+
                    \frac2{3m_{\rm path}}\right)E.
\]
The Rayleigh principle gives \eqref{eq:pathyes}. A path grounded
at both ends has energy $q(y)+y_1^2\geq a_\ell\|y\|^2$,
so the same lower bound holds for the whole matrix.

For the NO bound, choose an original uncovered edge $uv$.
Give $u,v$ value 1 and every other original vertex value 0.
On every length-$\ell$ path with endpoint values $x,y\in\{0,1\}$,
give its internal vertex $j$ the value
\begin{equation}\label{eq:pathtrial}
 z_j=\frac{x\sin((\ell-j)\theta_\ell)
               +y\sin(j\theta_\ell)}{\sin(\ell\theta_\ell)}.
\end{equation}
These values satisfy
$2z_j-z_{j-1}-z_{j+1}=b_\ell z_j$ at every internal vertex.
At each of $u,v$, there is one original path type with endpoint
values $(1,1)$ and two types with endpoint values $(1,0)$.
The total outward flux, divided by $m_{\rm path}$, is
\[
 3-\frac{3\sin((\ell-1)\theta_\ell)+\sin\theta_\ell}
             {\sin(\ell\theta_\ell)}=0,
\]
by \eqref{eq:pathroot} and the sine subtraction formula.
All other original vertices have value zero. If $D$ is the sum of
squared internal values across one copy of each original path type,
then the full trial vector has squared norm $MD+2$ and energy
exactly $b_\ell MD$.

All trial values lie in $[0,2]$. To see this, note that
$\ell\theta_\ell/2<\ell\pi/(4\ell-2)\leq\pi/3$ and
$\ell\theta_\ell<\pi$. The $(1,1)$ path values equal
\[
 \frac{\cos((j-\ell/2)\theta_\ell)}
      {\cos(\ell\theta_\ell/2)}\in[1,2],
\]
and the nonnegative $(1,0)$ and $(0,1)$ values are at most their
sum. In particular the $uv$ path alone gives $D\geq h$.

Now set the value of every internal vertex in $T$ to zero. The
original values already vanish on $T\cap V(H)$. Each altered
internal vertex has degree 2, so at most $2k$ edges change energy.
A changed edge has new energy at most 4, and the squared norm
falls by at most $4k$. The modified vector is nonzero because
$m_{\rm path} h\geq8k$ and it is an admissible Dirichlet test vector. Its
Rayleigh quotient is at most
\[
 \frac{b_\ell MD+8k}{MD+2-4k}
 \leq b_\ell+\frac{(8+4b_\ell)k}{MD-4k}
 \leq b_\ell+\frac{24k}{Mh},
\]
using $b_\ell<1$ and $MD\geq Mh\geq8k$.
This handles internal anchors without restricting the candidate set.

For completeness, the left side minus the right side of
\eqref{eq:pathroot} equals 2 at zero and is negative at
$\pi/(2\ell-1)$. Its derivative is negative in between, since
$\sin((\ell-\tfrac12)\theta)\geq\sin(\theta/2)>0$ there and
$3(\ell-\tfrac12)>1/2$. This proves existence, uniqueness, and
$0<b_\ell<a_\ell\leq1$. Put
$t_\ell=(\ell-\tfrac12)\theta_\ell\in(0,\pi/2)$.
Equation~\eqref{eq:pathroot} becomes
$3\cos t_\ell=\cos(t_\ell/(2\ell-1))$.
Its right side is strictly less than 1 and converges uniformly to 1,
so $t_\ell>\arccos(1/3)$ and
$t_\ell\to\arccos(1/3)$. Concavity of sine on $[0,\pi/2]$ gives
\[
 \frac{b_\ell}{a_\ell}
 =\left(\frac{\sin(t_\ell/(2\ell-1))}
 {\sin(\pi/(2(2\ell-1)))}\right)^2
 \geq\left(\frac{2t_\ell}{\pi}\right)^2>c_*.
\]
Using $2-2\cos x=4\sin^2(x/2)$
and $\sin x/x\to1$ proves \eqref{eq:pathlimit}.
\end{proof}

\begin{theorem}[Amplified exact-budget spectral hardness]
\label{thm:amplifiedhardness}
For every fixed $\gamma>c_*$, where $c_*$ is defined in
\eqref{eq:pathlimit}, a polynomial-time $\gamma$-approximation
for exact-budget grounded-$\lambda_{\min}$ maximization on connected
simple unweighted graphs would imply $\mathrm P=\mathrm{NP}$.
The reduction preserves the source budget $k$ exactly and permits
all graph vertices as candidates. In particular this excludes a
same-budget spectral factor $1-1/e$. It also extends to a prescribed
singleton seed and every fixed rational $\sigma\geq0$ on connected
positively weighted graphs. A randomized approximation with success
probability at least $2/3$ would imply $\mathrm{NP}=\mathrm{RP}$.
\end{theorem}
\begin{proof}
It suffices to consider $\gamma\leq1$. Choose a constant integer
$\ell\geq2$ with $b_\ell/a_\ell<\gamma$, possible by
\eqref{eq:pathlimit}, and put $\Delta=\gamma a_\ell-b_\ell>0$.
Fix an integer constant $C>\max\{8,32/\Delta\}$.
Given a connected cubic vertex-cover instance $(H,k)$,
$1\leq k<|V(H)|$, construct $G_{Ck,\ell}$ and keep budget $k$.
It has $n+Ckm(\ell-1)$ vertices and $Ckm\ell$ edges, so its size
is polynomial. Every vertex, including every path-interior vertex,
is eligible for grounding. Define the constants
\[
 \alpha=\frac{a_\ell}{1+4/(3C)},\qquad
 \beta=b_\ell+\frac{24}{C(\ell-1)}.
\]
Because $a_\ell\leq1$,
\[
 \gamma\alpha-\beta
 \geq\Delta-\frac{4\gamma a_\ell/3+24/(\ell-1)}C
 >\Delta-\frac{32}C>0.
\]
If $H$ has a cover of size at most $k$, enlarge it to size $k$
and ground those original vertices. Equation~\eqref{eq:pathyes}
gives value at least $\alpha$. If an arbitrary returned size-$k$
set $T$ has original part $T\cap V(H)$ that is not a cover,
\eqref{eq:pathno} gives value at most $\beta$.
Thus a $\gamma$-approximation on a YES instance must return an
original part that is a cover, since its value is at least
$\gamma\alpha>\beta$. Check that original part combinatorially.
On a NO instance no original part of size at most $k$ can pass.
This decides cubic vertex cover and proves the claimed hardness.
The constants depend only on the fixed target factor; strict
inequalities allow fixed rational upper/lower enclosures when
choosing them, independently of the input instance.

The special factor $1-1/e$ has an exact rational certificate.
For $\ell=4$, the ND path matrix and the sine recurrence give
$a_4$ as the smallest positive zero of $F$, and $3F(b_4)=1$, where
$F(x)=1-6x+5x^2-x^3$ is strictly decreasing on $[0,1/5]$.
Direct rational evaluations give
$F(99/500)>0$, $F(1/5)<0$, and $3F(309/2500)<1$.
Consequently
\[
 \frac{b_4}{a_4}<\frac{309/2500}{99/500}
 =\frac{103}{165}<\frac58<1-\frac1e.
\]
The last inequality follows from
$e>1+1+1/2+1/6+1/24=65/24>8/3$.

For a fixed seed and shift, multiply every path-edge weight by
a constant positive integer $W$, add a new mandatory seed $z$,
and join $z$ to every other vertex by a unit-weight edge.
After grounding $z$, any additional set $T$ has matrix
$W L(G_{Ck,\ell})_{-T}+(1+\sigma)I$.
Choose $W$ so that
\[
 W(\gamma\alpha-\beta)>(1-\gamma)(1+\sigma).
\]
The YES lower bound $W\alpha+1+\sigma$ and the NO upper bound
$W\beta+1+\sigma$ are then separated by factor $\gamma$.
The graph and its rational encoding remain polynomial, and the
budget still counts exactly $k$ additions to the mandatory seed.
The same combinatorial output check proves the reduction.
For randomized algorithms this check has no false positives and
succeeds with probability at least $2/3$ on YES instances, giving the $\mathrm{NP}=\mathrm{RP}$ consequence: an NP-complete
problem is decided with one-sided randomized error.
\end{proof}



\section{Marginal and Hybrid-State Identities}\label{app:maintenance}

For the scalar identities, $A\succ0$ is the current active matrix,
$b=e_v$, $A_u=A+ubb^\top$, and $R_t=(A+tI)^{-1}$.
Write $r_j(t)=b^\top R_t^j b$. The comparison matrix is
$B=A+\beta bb^\top$ for a finite penalty and $B=A[-v,-v]$ for
hard grounding. Thus $\Delta_f=\tr f(A)-\tr f(B)$.
\begin{lemma}[Scalar resolvent marginal]\label{lem:resolvent}
For $t\geq0$, $r_2(t)=-r'_1(t)$ and
\begin{equation}\label{eq:qone}
 q_1(t):=\tr(A+tI)^{-1}-\tr(B+tI)^{-1}
 =\begin{cases}
 \displaystyle\frac{\beta r_2(t)}{1+\beta r_1(t)},&\beta<\infty,\\[6pt]
 \displaystyle\frac{r_2(t)}{r_1(t)},&\beta=\infty.
 \end{cases}
\end{equation}
Consequently, for $0<s<1$,
\begin{equation}\label{eq:fractionaloracle}
 \Delta_s=\frac{\sin(\pi s)}{\pi}
              \int_0^\infty t^{-s}q_1(t)dt,
 \qquad \Delta_1=q_1(0).
\end{equation}
\end{lemma}
\begin{proof}
Differentiating $(A+tI)R_t=I$ gives $R'_t=-R_t^2$, hence
$r'_1=-r_2$. For the finite update, direct multiplication verifies
the Sherman--Morrison identity
\[
 (A+tI+\beta bb^\top)^{-1}
 =R_t-\frac{\beta R_tbb^\top R_t}{1+\beta b^\top R_tb}.
\]
Taking traces gives the first formula. For hard grounding reorder
the candidate first and write $R_t$ in blocks on $\{v\},U$. The
block inverse identity gives
\[
 (C+tI)^{-1}=(R_t)_{UU}
       -\frac{(R_t)_{Uv}(R_t)_{vU}}{(R_t)_{vv}}.
\]
Its trace difference from $R_t$ is
$r_1+\|(R_t)_{Uv}\|^2/r_1=r_2/r_1$; the formula also holds when
$U$ is empty. The scalar identity
\[
 x^{-s}=\frac{\sin(\pi s)}{\pi}
           \int_0^\infty\frac{t^{-s}}{x+t}dt
\]
follows from $t=xu$ and the beta integral
$\int_0^\infty u^{-s}/(1+u)du=\Gamma(s)\Gamma(1-s)=\pi/\sin(\pi s)$.
Diagonalization gives the corresponding trace identity. The integrals
for each trace are finite, so subtracting them and using
\eqref{eq:qone} proves \eqref{eq:fractionaloracle}. At $s=1$ use
the trace difference at $t=0$ directly.
\end{proof}

\begin{proposition}[Exact two-moment updates]\label{prop:updates}
At one shift let $R=(A+tI)^{-1}$, $x=Re_v$, $y=R^2e_v$, and
$s_x=r_{2,v}=x^\top x$. Put $c=(\beta^{-1}+r_{1,v})^{-1}$, using
$\beta^{-1}=0$ for hard grounding. Every remaining candidate $u$ obeys
\begin{align}
 r_{1,u}^+&=r_{1,u}-c x_u^2,\label{eq:updateone}\\
 r_{2,u}^+&=r_{2,u}-2c x_u y_u+c^2 s_x x_u^2.\label{eq:updatetwo}
\end{align}
Once $x,y$ are available, updating all candidates at $q$ shifts costs
$O(nq)$ arithmetic and $O(nq)$ cache storage. Selected-column generation
requires two solves per shift, or one retained Krylov basis from $e_v$
with both small functions $(T+tI)^{-1}e_1,(T+tI)^{-2}e_1$.
\end{proposition}
\begin{proof}
Sherman--Morrison gives $R^+=R-cxx^\top$ for soft grounding.
For hard grounding the same matrix has zero row and column $v$;
its complementary block is $(A[-v,-v]+tI)^{-1}$ by the block inverse
identity in Lemma~\ref{lem:resolvent}. Thus diagonal entries, and
squared-column norms at remaining coordinates, can be computed before
or after deleting the zero row and column. Squaring gives
\[
 (R-cxx^\top)^2=R^2-c(yx^\top+xy^\top)+c^2(x^\top x)xx^\top.
\]
Taking diagonals proves both updates. Two shifted solves produce
$x$ and $y$; a retained Lanczos basis gives their projected versions
from the two indicated small vectors, exactly at termination.
Reconstructing both vectors at all shifts costs $O(nqh)$ after an
$h$-step basis is stored, in addition to the large matrix--vector
products. Only the cache update itself costs $O(nq)$.
\end{proof}


For the following all-order formula, $R$ is the current zero-padded
hard-grounded inverse, $x=Re_v$, and $c=1/x_v$.
The notation $S$ includes the fixed seed, and $\bar R_v=R-cxx^\top$.
\begin{proposition}[All-order independent decrements]\label{prop:incall}
Let $R=X_S^{-1}$ be zero padded on already selected coordinates,
let $v$ be active, $x=Re_v$, $c=1/R_{vv}$, and
$\bar R_v=R-cxx^\top$. For $i\geq0$ define $a_i=R^ix$,
$s_i=x^\top R^ix$, and, for each active $u$,
\[
 \gamma_0(u)=a_0(u),\qquad
 \gamma_t(u)=a_t(u)-c\sum_{i=0}^{t-1}s_i\gamma_{t-1-i}(u).
\]
Then $\gamma_t(u)=e_u^\top \bar R_v^tx\geq0$ and the moment decrement is
\[
 D_j(u)=r_j(u,S)-r_j(u,S\cup\{v\})
      =c\sum_{i=0}^{j-1}a_i(u)\gamma_{j-1-i}(u)\geq0.
\]
These decrements depend on selected columns, not on the cached
candidate moments. After a last refresh at time $b$,
\[
 |\widehat r_j(u,S_t)-r_j(u,S_t)|
 \leq E_{j,u}^{(b)}+\sum_{i=b}^{t-1}\varepsilon_{j,u}^{(i)},
 \qquad
 |\widehat D_j^{(i)}-D_j^{(i)}|\leq\varepsilon_{j,u}^{(i)}.
\]
For exact columns all $\varepsilon$ are zero. In an implementation
whose refresh flag changes only from false to true, whose refresh
batch contains only false-flag candidates, and whose nonseed
selections always have true flags, a run starting with all $N$
candidate flags false has exactly $N-n_{\rm ref}$ never-refreshed candidates
after $n_{\rm ref}$ total refreshes. All of them remain unselected.
\end{proposition}
\begin{proof}
For small $|z_{\rm ser}|$, the rank-one inverse identity gives
\[
 (I-z_{\rm ser}\bar R_v)^{-1}x=
 \frac{(I-z_{\rm ser}R)^{-1}x}{1+c z_{\rm ser}x^\top(I-z_{\rm ser}R)^{-1}x}.
\]
Equating coefficients proves the recurrence for $\bar R_v^tx$.
The telescoping identity
\[
 R^j-\bar R_v^j=\sum_{i=0}^{j-1}R^i(R-\bar R_v)\bar R_v^{j-1-i}
\]
proves the displayed decrement. Grounded inverses and their zero-padded
versions are entrywise nonnegative, so every displayed product is
nonnegative. Subtracting the computed update from the true update
and telescoping proves the error envelope; refreshing restarts it
with that refresh's error budget. Finally each refresh changes one
distinct false flag to true, so exactly $n_{\rm ref}$ flags have changed.
Only true flags can be selected, proving the last assertion.
\end{proof}


We next give the maintained-solution sketch result with $A_0=C^\top C$,
$W=V\setminus S_0$, $N=|W|$, and $R_T$ the zero-padded inverse.
\begin{proposition}[Current-matrix sketch maintenance for hard CFCC]
\label{prop:maintainedsketch}
Let $A_0=X_{S_0}=C^\top C\succ0$ have dimension $N$.
For $T\subset W$, write $R_T$ for the zero-padded inverse of
$A_0[W\setminus T,W\setminus T]$. Fix matrices $\Pi_C,\Pi_I$ and maintain
\[
 Y_1=\Pi_C C R_T,\qquad Y_2=\Pi_I R_T.
\]
When grounding an active vertex $v$, one solve for $x=R_Te_v$
updates both sketches exactly by
\[
 Y_i' =Y_i-\frac{Y_i[:,v]x^\top}{x_v},\qquad i=1,2.
\]
For every fixed $T$, if $\Pi_C,\Pi_I$ are independent Gaussian matrices
with $d$ rows and entry variance $1/d$, then
$\widehat r_1(u)=\|Y_1[:,u]\|^2$ and
$\widehat r_2(u)=\|Y_2[:,u]\|^2$ are unbiased estimates of the
\emph{current} moments $r_1(u,T),r_2(u,T)$.

For $0<\zeta<1$, $0<\delta<1$, and a maximum of $k<N$ new anchors,
put $\mathcal N_{N,k}=\sum_{t=0}^k\binom Nt(N-t)$. If
\[
 d\geq8\zeta^{-2}\log(4\mathcal N_{N,k}/\delta),
\]
both estimates have simultaneous relative error at most $\zeta$
at every active vertex of every such set, with probability at
least $1-\delta$. This event covers adaptive choices made from
the maintained sketches. Choosing the maximum estimated ratio
$\widehat r_2/\widehat r_1$, or exactly rescoring a candidate
pool containing that maximizer and choosing its best member,
achieves a maximum-marginal fraction at least
$\alpha_{\rm step}=((1-\zeta)/(1+\zeta))^2$. For $k$ selections its
normalized gain is at least $(1-e^{-\alpha_{\rm step}})\mathrm{OPT}_k$.
In particular $\zeta\leq\epsilon/4$ gives the
$(1-e^{-1}-\epsilon)$ factor. The sufficient row count has order
$O(\zeta^{-2}((k+1)\log(N+1)+\log(1/\delta)))$.

Each update costs $O(Nd)$ scalar work and uses $O(Nd)$ sketch
storage. The selected-column solve can be reused from rescoring;
no $R_T^2e_v$ solve is needed for this sketch update.

Moreover, for any row-vector error $e$ supported on the active
coordinates, its exact-column update $e'=e-e_vx^\top/x_v$ obeys
\[
 \|e'\|_{A_0}^2=\|e\|_{A_0}^2-e_v^2/x_v.
\]
Thus accumulated initial sketch-solve errors are contractive in
the fixed $A_0$ energy norm under exact subsequent columns.
\end{proposition}
\begin{proof}
Sherman--Morrison in its hard zero-padded form gives
$R_{T+v}=R_T-xx^\top/x_v$. Multiplying on the left by $\Pi_C C$ or
$\Pi_I$ proves the update because $Y_i[:,v]$ is that left factor
times $x$. On active coordinates $R_TA_0R_T=R_T$, and all
remaining rows and columns of both sides are zero. Consequently
\[
 \|CR_Te_u\|^2=r_1(u,T),\qquad
 \|R_Te_u\|^2=r_2(u,T).
\]
For a fixed deterministic set and candidate, Gaussian norm
concentration gives failure at most $2e^{-d\zeta^2/8}$ for each
of the two vectors. There are $2\mathcal N_{N,k}$ vectors across all sets and
candidates, so the stated row count gives the joint event.
Since it holds simultaneously for this deterministic family,
it also holds along every adaptive selection path. The bound
$\mathcal N_{N,k}\leq N(k+1)N^k$ gives the stated asymptotic dimension.

On the event each ratio is between $(1-\zeta)/(1+\zeta)$ and
$(1+\zeta)/(1-\zeta)$ times its true marginal. Comparing an
estimated maximizer to a true maximizer therefore gives $\alpha_{\rm step}$.
Exact rescoring a pool containing that candidate can only improve
the true marginal at the same current set. The the integer analysis induction gives
$1-e^{-\alpha_{\rm step}}$; since $1-\alpha_{\rm step}=4\zeta/(1+\zeta)^2\leq4\zeta$,
the final approximation follows. The two rank-one outer products
and recomputation of column norms use $O(Nd)$ operations and
storage, in addition to one selected-column solve.

For a supported row $e$, the active inverse equations imply
$eA_0x=e_v$ and $x^\top A_0x=x_v$. Expanding the squared energy
norm gives the last identity. The new row has zero at $v$ and
at every previous anchor, so the contraction can be iterated.
\end{proof}

\paragraph{Initial solve errors in the maintained sketches.}
For a matrix error $\mathcal E$ with supported rows, define
$\|\mathcal E\|_{F,A_0}^2=\tr(\mathcal E A_0\mathcal E^\top)$.
Cauchy--Schwarz with the representer $R_Te_u$ gives
$\|\mathcal E[:,u]\|\leq\sqrt{r_1(u,T)}\|\mathcal E\|_{F,A_0}$.
Let $B\geq\lambda_{\max}(A_0)$ and let $e_i$ bound the initial
energy errors for the two solution sketches. The contraction above
keeps those bounds valid under exact-column updates. On the
simultaneous Gaussian event, the relative squared-norm errors
are at most
\[
 \zeta+2\sqrt{1+\zeta}\,e_1+e_1^2,
 \qquad
 \zeta+2\sqrt{1+\zeta}\sqrt B\,e_2+B e_2^2,
\]
respectively. To obtain the second bound, use the spectral
inequality $r_1\leq B r_2$. Both follow by expanding the
perturbed norm square and bounding the cross term. 



\section{Adaptive Concentration and Capped Greedy}\label{app:adaptive}

The first statement covers $J=p+1$ moments, and specializes to the
two CFCC solution sketches for $p=1$. The second gives the stopping-rule
analysis for both general integer scores and the sharper CFCC ratio.
\begin{theorem}[Uniform adaptive moment accuracy]
\label{thm:adaptivehybridjl}
Let all entries of $\Pi_C,\Pi_I$ be independent centered Gaussians of
variance $1/d$, or independent signs $\pm1/\sqrt d$.
For $J=p+1$, $0<\zeta,\delta<1$, put
\[
 \mathcal N_{N,k}=\sum_{t=0}^k\binom Nt(N-t),\qquad
 d_{\rm reuse}=\left\lceil
    8\zeta^{-2}\log\frac{2J\mathcal N_{N,k}}{\delta}\right\rceil.
\]
With $d\geq d_{\rm reuse}$ and exact sketch maintenance,
all $J$ moment estimates have relative error at most $\zeta$
at every active vertex of every $|T|\leq k$, with probability
at least $1-\delta$. The same event covers every adaptive
selection rule, refresh pattern, and stopping decision.
With $\zeta\leq\eta/(16p)$, the computed integer marginals
have relative error at most $\eta$ by
Theorem~\ref{thm:polymoment}.

For comparison, new independent sketches in each of $k$ rounds
need only
\[
 d_{\rm fresh}=\left\lceil
       8\zeta^{-2}\log\frac{2JNk}{\delta}\right\rceil
\]
for the same ideal-solve moment event along the chosen path.
In particular, the sufficient dimensions have orders
\begin{align*}
 d_{\rm reuse}
 &=O\!\left(p^2\eta^{-2}
       \left[(k+1)\log(N+1)+\log(p/\delta)\right]\right),\\
 d_{\rm fresh}
 &=O\!\left(p^2\eta^{-2}\log(Npk/\delta)\right).
\end{align*}
\end{theorem}
\begin{proof}
For a deterministic nonzero vector $w$, let $Z$ be one
unscaled sketch row dotted with $w/\|w\|$. In either distribution,
$\E Z^2=1$, $\E Z^4\leq3$, and
$\E e^{tZ}\leq e^{t^2/2}$. For signs the last inequality follows
from $\cosh x\leq e^{x^2/2}$ and independence. Integrating
against an auxiliary standard Gaussian yields
$\E e^{tZ^2}\leq(1-2t)^{-1/2}$ for $0<t<1/2$.
Chernoff's bound with $t=\zeta/[2(1+\zeta)]$ bounds the upper
tail of the average of $d$ independent $Z^2$ by
\[
 \exp\{-d[\zeta-\log(1+\zeta)]/2\}
 \leq e^{-d\zeta^2/8}.
\]
Here $\zeta-\log(1+\zeta)\geq\zeta^2/4$ on $[0,1]$, as
follows by differentiating the difference.
For the lower tail, $e^{-y}\leq1-y+y^2/2$ for $y\geq0$
gives $\E e^{-tZ^2}\leq\exp(-t+3t^2/2)$.
Taking $t=\zeta/3$ gives $e^{-d\zeta^2/6}$.
The combined two-sided failure is at most $2e^{-d\zeta^2/8}$.

The targets $R_T^\ell e_u$ and $CR_T^\ell e_u$ form a
\emph{deterministic} family with $J\mathcal N_{N,k}$ members when only
the first $J$ moments are retained. They are nonzero at active
vertices. The preceding tail bound and a union bound prove the
claim. Their correlations, including shared rows across powers,
do not affect this union bound. Once the event holds, adaptively
choosing a member of this family cannot invalidate it.
The higher-power identity computes precisely these fixed-randomness
sketches. Unbiasedness for each deterministic $T$ supplies the
moment identity, while this simultaneous event supplies the
adaptive guarantee.

With fresh randomness, condition on the complete history at each
round; its current matrix is then fixed independently of the new
rows. A union bound over at most $NJ$ current moments and a
conditional failure budget $\delta/k$ proves the fresh formula.
Finally $\mathcal N_{N,k}\leq N(k+1)N^k$ and
$\zeta=\eta/(16p)$ prove the stated orders.
\end{proof}

\begin{theorem}[Hybrid retention and capped refreshes]
\label{thm:hybridcap}
At each state, suppose every not-yet-refreshed marginal estimate
satisfies
$(1-\eta)\Delta_u\leq\widehat\Delta_u\leq(1+\eta)\Delta_u$,
and refreshed candidates are maintained exactly by the incremental analysis.
Use the estimate itself as the heap key for an unrefreshed
candidate and the exact marginal for a refreshed one.
Let $0\leq\theta<1$ be an acceptance tolerance. Process keys
in descending order, exactly rescoring batches of unrefreshed
candidates and retaining the best exact value $b$. Stop either
when $b\geq(1-\theta)q_{\rm heap}$, where $q_{\rm heap}$ is the largest remaining
key, or after a positive refresh cap has been reached, and
return the best processed candidate. Every processed refreshed
candidate is included in this comparison. If a budget stop occurs,
at least the initial maximum-key candidate has been processed.
Then the selected marginal is at least
\[
 \alpha_{\rm step}\max_u\Delta_u,\qquad
 \alpha_{\rm step}=\min\left\{\frac{1-\eta}{1+\eta},
                           (1-\theta)(1-\eta)\right\}.
\]
The $k$-round normalized gain is at least
$(1-e^{-\alpha_{\rm step}})\mathrm{OPT}_k$.
In particular $\eta\leq\epsilon/4$ and $\theta\leq\epsilon/2$
give $(1-1/e-\epsilon)\mathrm{OPT}_k$ on the uniform event of
Theorem~\ref{thm:adaptivehybridjl}.

For $p=1$, put $\chi=(1+\zeta)/(1-\zeta)$ on its moment event.
The sharper factor is
$\alpha_{\rm step}=\min\{\chi^{-2},(1-\theta)\chi^{-1}\}$; thus
$\zeta\leq\epsilon/8$ and $\theta\leq\epsilon/2$ suffice.
Arbitrarily many candidates can be kept exact after their first
refresh. Their total number of refreshes is $R\leq N$.
\end{theorem}
\begin{proof}
Both types of key lie between $(1-\eta)\Delta_u$ and
$(1+\eta)\Delta_u$. If $u_*$ maximizes the true marginal and
$w$ has maximum initial key, then
\[
 \Delta_w\geq\frac{\mathrm{key}(w)}{1+\eta}
 \geq\frac{1-\eta}{1+\eta}\Delta_{u_*}.
\]
A budget stop has processed $w$ and returns a candidate no worse
than $w$. At an acceptance stop, every processed true value is
at most $b$, and every unprocessed one is at most $q_{\rm heap}/(1-\eta)$.
Thus $b\geq(1-\theta)(1-\eta)\Delta_{u_*}$.
Taking the minimum covers both stopping rules, including an
empty remaining heap. Positive-budget batching may overshoot
its nominal size without changing this argument.

By Theorem~\ref{thm:cm} the gain is normalized, monotone, and submodular relative
to $S_0$. At each round its remaining gap to the best $k$-element
extension is at most $k$ times the largest current marginal.
The gap therefore contracts by at least $1-\alpha_{\rm step}/k$ at every
step, proving $1-(1-\alpha_{\rm step}/k)^k\geq1-e^{-\alpha_{\rm step}}$.
Moreover $1-\alpha_{\rm step}\leq\max\{2\eta,\eta+\theta\}\leq\epsilon$,
and $e^{-\alpha_{\rm step}}-e^{-1}\leq1-\alpha_{\rm step}$ proves the claimed factor.
For $p=1$ each ratio key is between $\Delta_u/\chi$ and
$\chi\Delta_u$. Repeat the same proof with these endpoints.
Here $1-\chi^{-2}\leq4\zeta$ and
$1-(1-\theta)\chi^{-1}\leq\theta+2\zeta$.
Exact retention follows from the independent-column decrement
identity in the incremental analysis; each false-to-true refresh flag changes at most
once. Replacing any sketched value by its exact value preserves
all the preceding key inequalities, even for an adaptively
chosen subset of candidates.
\end{proof}



\section{Numerical Error Budgets}\label{app:numerical}

In the following proposition the spectral bounds $a,B$ refer to
$A_0$, and the Frobenius energy norm is
$\|\mathcal E\|_{F,A_0}^2=\tr(\mathcal E A_0\mathcal E^\top)$. Zero padding is enforced on
all grounded coordinates. The final paragraph applies to retained
scalar entries and uses the independent decrements from
Proposition~\ref{prop:incall}.
\begin{proposition}[Residual budgets for maintained sketches]
\label{prop:sketchresidual}
Let $a\leq\lambda_{\min}(A_0)$ and
$B\geq\lambda_{\max}(A_0)$. At a $p=1$ update, let
$\widetilde x$ be supported on the current active set $U$, with
$\widetilde x_v>0$ and
$\|A_0[U,U]\widetilde x_U-e_v\|_2\leq\rho_{\rm res}$.
For either sketch, perform
\[
 \widetilde Y'=
 \widetilde Y-\widetilde Y[:,v]\widetilde x^\top/
                 \widetilde x_v+D,
\]
where $D$ has zero columns on all new and old anchors and
$\|D\|_{F,A_0}\leq\omega$.
If $e\geq\|\widetilde Y-Y\|_{F,A_0}$ before the update, then
\begin{equation}\label{eq:sketchresidual}
 e'=e+\frac{\|\widetilde Y[:,v]\|_2\rho_{\rm res}}
                   {\sqrt a\,\widetilde x_v}+\omega
\end{equation}
is a valid bound afterwards. This recurrence is deterministic
along adaptive selections. At every state the two moment errors
are bounded by
\[
 \zeta+2\sqrt{1+\zeta}\,e_1+e_1^2,\qquad
 \zeta+2\sqrt{1+\zeta}\sqrt B\,e_2+B e_2^2
\]
on the ideal-sketch event, where $e_1,e_2$ are their current
accumulated energy bounds.

For higher powers, fix a left factor $\mathcal F$. A sufficient
alternative certificate uses active-coordinate residuals
$\rho_\ell=\|\widetilde{\mathcal Y}_\ell A_U-\widetilde{\mathcal Y}_{\ell-1}\|_F$,
with $\widetilde{\mathcal Y}_0=\mathcal F[:,U]$ exact. Recursively let
$E_0=0$ and $E_\ell=(E_{\ell-1}+\rho_\ell)/a$.
Then $\|\widetilde{\mathcal Y}_\ell-\mathcal F[:,U]A_U^{-\ell}\|_F\leq E_\ell$.
For even moments take $e=B^\ell E_\ell$, and for odd moments
take $e=B^{\ell-1/2}E_\ell$. Their relative squared-norm error
is at most $\zeta+2\sqrt{1+\zeta}\,e+e^2$.
\end{proposition}
\begin{proof}
Let $x=R_Te_v$, $h=\widetilde x-x$, and $h_v=\widetilde x_v-x_v$.
Active inverse equations give
$\|h\|_{A_0}\leq\rho_{\rm res}/\sqrt a$,
$\langle h,x\rangle_{A_0}=h_v$, and
$\|x\|_{A_0}^2=x_v$. Therefore
\[
 \left\|\frac{\widetilde x}{\widetilde x_v}
              -\frac{x}{x_v}\right\|_{A_0}^2
 =\frac{\|h\|_{A_0}^2-h_v^2/x_v}{\widetilde x_v^2}
 \leq\frac{\rho_{\rm res}^2}{a\widetilde x_v^2}.
\]
Write $P=I-e_vx^\top/x_v$ and $E=\widetilde Y-Y$.
Subtracting the exact update yields
\[
 \widetilde Y'-Y'=EP+
 \widetilde Y[:,v]\left(x/x_v-\widetilde x/\widetilde x_v\right)^\top+D.
\]
Proposition~\ref{prop:maintainedsketch} contracts $EP$ in energy.
The energy norm of an outer product is the product of the
Euclidean norm of its left vector and the energy norm of its
right vector. The triangle inequality proves
\eqref{eq:sketchresidual}. All new errors remain supported on
unselected coordinates, so the argument iterates. Applying the
coordinate-representer bounds after
Proposition~\ref{prop:maintainedsketch} proves the moment bounds.

For higher powers, subtract the exact chain equation at each
level and multiply by $A_U^{-1}$, whose norm is at most $1/a$
by principal-submatrix interlacing. This proves the $E_\ell$
recurrence. The true target squared norms are at least
$B^{-2\ell}$ and $B^{-(2\ell-1)}$, respectively. On the ideal
JL event, their sketched norm is at most $\sqrt{1+\zeta}$
times the true target norm. Expanding a perturbed squared norm
and bounding its cross term proves the final error expression.
\end{proof}

\paragraph{Instantiating numerical budgets.}
For example, a $p=1$ target moment tolerance $\bar\zeta<1$
is met by an ideal-sketch tolerance $\zeta=\bar\zeta/2$ and
$e_1,\sqrt B e_2\leq\bar\zeta/8$; substitution in the preceding
bound is at most $\bar\zeta$. Initial solve errors, subsequent
column residuals, and rounding terms share these budgets.
The same substitution applies to the higher-power normalized
errors $e$. Refreshed scalar entries use the the incremental analysis decrement
intervals and refresh-error bounds; if their resulting relative
marginal error is at most $\nu$, replacing exact values throughout
the hybrid proof loses at most the additional factor
$(1-\nu)/(1+\nu)$, with key tolerance enlarged to
$\max\{\eta,\nu\}$. This follows by comparing the largest
computed processed value with the largest true processed value
before applying the two stopping arguments. Residual evaluation
for every higher-power level costs $O(dL(m+N))$ per checked
state and is charged separately when used. These statements
specify numerical conditions, in addition to the sketch event,
for an inexact implementation of the preceding algorithms.

\paragraph{Scalar decrement errors}
Let a selected-column chain have true-residual bounds $\rho_i$ and
spectral lower bound $a_*>0$. With $a_{-1}=e_v$, its Euclidean
errors satisfy $E_{-1}=0$, $E_i\leq(E_{i-1}+\rho_i)/a_*$.
Use component intervals $\widehat a_i(u)\pm E_i$,
and intervals for $s_i$ centered at $\widehat a_0^\top\widehat a_i$
with radius
$\|\widehat a_0\|E_i+\|\widehat a_i\|E_0+E_0E_i$.
The reciprocal of a positive interval for $a_0(v)$ encloses $c$.
Propagating these intervals through the finite $\gamma_t,D_j$
recurrences encloses every true decrement by induction. Their
distance from the computed decrement gives the $\varepsilon$
above. This construction has the same $O(NJ^2)$ scalar order as
the update. The residual inequality follows by applying $R$ to
each chain equation; the inner-product radius follows by expansion
and Cauchy--Schwarz. 



\section{Seed-Preconditioner Reuse}\label{app:preconditioning}
Write $W=V\setminus S_0$, $N=|W|$, $A_0=X_{S_0}\succ0$,
$T=S\setminus S_0$, $t=|T|$, and $U=W\setminus T$.
All inverses in this subsection act on their active coordinates.
An inverse preconditioner $B_0\succ0$ for $A_0$ is applied to a
vector on $U$ by zero extension to $W$, application of $B_0$, and
restriction back to $U$. Its active matrix is $B_U=B_0[U,U]$.

\begin{theorem}[Masked seed preconditioning]\label{thm:seedpcg}
Suppose $aA_0^{-1}\preceq B_0\preceq bA_0^{-1}$ for $0<a\leq b$.
Let $A=A_0[U,U]$ and $D_U=|U|>0$. Every eigenvalue of
$P_U=A^{1/2}B_UA^{1/2}$ is at least $a$, and at most $t$ of its
eigenvalues exceed $b$. Let $n_{\rm out}\leq\min\{t,D_U\}$ count those
exceeding $b$, and set $\chi_c=b/a$. In exact arithmetic PCG obeys,
for $h\geq n_{\rm out}$ and $\chi_c>1$,
\[
 \frac{\|e_h\|_A}{\|e_0\|_A}
 \leq 2\left(\frac{\sqrt{\chi_c}-1}{\sqrt{\chi_c}+1}\right)^{h-n_{\rm out}}.
\]
It therefore reaches relative energy error $\xi$ in at most
\[
 \min\left\{D_U,\ t+
 \left\lceil\frac{\log(2/\xi)}{
 \log((\sqrt{\chi_c}+1)/(\sqrt{\chi_c}-1))}\right\rceil\right\}
\]
iterations. For $a=b$, termination occurs within $\min\{D_U,t+1\}$
iterations. These bounds hold simultaneously for every $T$ on the
single seed spectral-approximation event.

The identity-padded matrix $\overline A=A\oplus I_T$ differs from
$A_0$ by rank at most $2t$.
\end{theorem}
\begin{proof}
Order $U$ before $T$ and write
$A_0=\left(\begin{smallmatrix}A&C\\C^\top&D\end{smallmatrix}\right)$.
The block inverse formula gives
\[
 (A_0^{-1})[U,U]
 =A^{-1}+A^{-1}C(D-C^\top A^{-1}C)^{-1}C^\top A^{-1}.
\]
The Schur complement is positive definite. Congruence by $A^{1/2}$
thus yields $\mathcal K=I+ZZ^\top$ with $\operatorname{rank}Z\leq t$.
Compression preserves the assumed Loewner inequalities, so
$a\mathcal K\preceq P_U\preceq b\mathcal K$. In particular $P_U\succeq aI$.
On the subspace $\ker Z^\top$, of codimension at most $t$,
the Rayleigh quotient of $P_U$ is at most $b$. The min--max
principle proves the outlier count. The preconditioned system
$B_U^{1/2}AB_U^{1/2}$ has the same spectrum as $P_U$.

For completeness, CG minimizes the energy norm of the error over
its affine Krylov space: its conjugate search directions are
pairwise orthogonal in the system inner product and exact line
minimization makes the final error orthogonal to their span.
Apply this property to $P_{\rm pcg}=B_U^{1/2}AB_U^{1/2}$ and the transformed
error $\widetilde e=B_U^{-1/2}e$, for which
$\|\widetilde e\|_{P_{\rm pcg}}=\|e\|_A$. The Krylov error candidates are
$P(P_{\rm pcg})\widetilde e_0$ for polynomials of degree at most $h$ with
$P(0)=1$. Diagonalizing $P_{\rm pcg}$ bounds their $P_{\rm pcg}$-norm by
$\max_{\lambda\in\spec(P_{\rm pcg})}|P(\lambda)|\|e_0\|_A$.
Since $P_{\rm pcg}$ and $P_U$ have the same spectrum, this is the required
polynomial bound in terms of $P_U$.

Let $\mu_1,\ldots,\mu_{n_{\rm out}}>b$ be the outliers, with multiplicity.
Multiply the degree-$(h-n_{\rm out})$ Chebyshev residual polynomial on
$[a,b]$, normalized to equal one at zero, by
$\prod_{i=1}^{n_{\rm out}}(1-z/\mu_i)$. The product vanishes at every outlier
and has absolute value at most one on $[a,b]$. The Chebyshev
factor is bounded by
$2((\sqrt{\chi_c}-1)/(\sqrt{\chi_c}+1))^{h-n_{\rm out}}$ there, proving the result.
For $a=b$ the polynomial $(1-z/a)\prod_i(1-z/\mu_i)$ annihilates
the spectrum. The usual degree-$D_U$ annihilating polynomial gives
the dimension cap. Finally
\[
 \overline A-A_0=
 \begin{pmatrix}0&-C\\-C^\top&I-D\end{pmatrix}
 =\begin{pmatrix}0&-C\\0&I-D\end{pmatrix}
  +\begin{pmatrix}0&0\\-C^\top&0\end{pmatrix},
\]
a sum of two matrices of rank at most $t$. The stronger one-sided
outlier conclusion follows from the block inverse calculation,
rather than from this rank count alone. Every argument is
deterministic for every principal subset once $B_0$ satisfies its
initial spectral inequalities.
\end{proof}

\paragraph{Approximate Cholesky instantiation.}
Contract the seed to a vertex $g$, and represent any diagonal
killing by edges to $g$. Let $\mathcal L$ be the resulting connected
Laplacian, whose principal matrix off $g$ is $A_0$. The
SparseCholesky theorem of Kyng--Sachdeva~\cite{kyng2016},
Theorem 3.1, supplies $\widetilde{\mathcal L}$ with
$(1-\epsilon_c)\mathcal L\preceq\widetilde{\mathcal L}
\preceq(1+\epsilon_c)\mathcal L$ except with probability $\delta_c$.
Its expected construction work and factor storage are
$O(\epsilon_c^{-2}(m_E+N)\log^2((N+1)/\delta_c)\log(N+1))$;
Appendix B gives the high-probability work variant with the stated
sampling rule. The sparse inverse application has work $O(F+N)$,
where $F$ is the stored factor size.

Define $Q$ to have rows $e_u^\top-e_g^\top$, $u\in W$.
Then $B_0=Q\widetilde{\mathcal L}^{+}Q^\top$ equals the inverse
of $\widetilde{\mathcal L}[W,W]$: solving with balanced source
$(r,-\ind^\top r)$ and subtracting the ground potential solves
the grounded system. Restriction of the Laplacian spectral
inequalities and inversion give
$a=(1+\epsilon_c)^{-1}$, $b=(1-\epsilon_c)^{-1}$, hence
$\chi_c=(1+\epsilon_c)/(1-\epsilon_c)$. At $\epsilon_c=1/2$,
$\chi_c=3$ and the iteration bound is $t+O(\log(1/\xi))$.
This instantiation uses the imported factorization theorem;
Theorem~\ref{thm:seedpcg} proves its simultaneous reuse under
grounding. 



\section{Integer Marginals and Higher-Power Sketches}\label{app:integer}

These results prove Proposition~\ref{main:integer} and develop
the higher-power maintained-sketch extension. In the first two
statements, $A\succ0$, $b=e_v$, $r_j=b^\top(A+tI)^{-j}b$, and
$\beta^{-1}=0$ denotes the hard case. The logarithmic recurrence
is given explicitly in Proposition~\ref{prop:allpowers} below.
\begin{proposition}[Positive integer marginal representation]\label{prop:integerpositive}
At any nonnegative shift let $R=(A+tI)^{-1}$, $x=Rb$,
$c=(\beta^{-1}+r_1)^{-1}$, using $\beta^{-1}=0$ for hard grounding,
and $\bar R_v=R-cxx^\top\succeq0$. Define
\begin{gather*}
 a_j=cr_{j+1},\qquad d_j=cx^\top \bar R_v^{j-1}x\quad(j\geq1),\\
 H_0=1,\qquad H_j=\sum_{i=1}^j d_i H_{j-i}.
\end{gather*}
The integer marginal admits the noncancelling formula
\begin{equation}\label{eq:integerpositive}
 q_p(t)=\sum_{j=1}^p j d_j H_{p-j},\qquad
 a_p\leq q_p(t)\leq p a_p.
\end{equation}
All terms are nonnegative. In particular, with $\nu$ the probability
measure defined below,
\begin{equation}\label{eq:q2variance}
 q_2=2a_2-a_1^2=a_2+\operatorname{Var}_{\nu}(Z),
 \qquad a_1^2\leq a_2.
\end{equation}
\end{proposition}
\begin{proof}
For soft grounding $\bar R_v=(A+tI+\beta bb^\top)^{-1}$.
For hard grounding it is the zero-padded inverse of the deleted
principal matrix. These facts follow from the principal-inverse identities,
so $q_p=\tr R^p-\tr \bar R_v^p$ and $\bar R_v\succeq0$ in both cases.
Write $D(z_{\rm ser})=\sum_{j\geq1}d_jz_{\rm ser}^j=c z_{\rm ser}x^\top(I-z_{\rm ser}\bar R_v)^{-1}x$.
For sufficiently small $|z_{\rm ser}|$, the determinant lemma gives
\begin{gather*}
 \frac{\det(I-z_{\rm ser}\bar R_v)}{\det(I-z_{\rm ser}R)}=\frac1{1-D(z_{\rm ser})},\\
 \sum_{p\geq1}\frac{q_p}{p}z_{\rm ser}^p=-\log(1-D(z_{\rm ser})).
\end{gather*}
Every $d_j$ is nonnegative by positive semidefiniteness of $\bar R_v$.
The reciprocal power series has coefficients $H_j$ satisfying the
stated positive recurrence. Differentiating the logarithmic series
proves \eqref{eq:integerpositive}'s equality.

Alternatively expand words in $R=\bar R_v+cxx^\top$ inside $\tr R^p$.
Marking one of the $\ell$ occurrences of $cxx^\top$ and cyclically
rotating the trace gives
\[
 q_p=\sum_{\ell=1}^p\frac p\ell
       \sum_{\substack{j_1+\cdots+j_\ell=p\\j_i\geq1}}
       d_{j_1}\cdots d_{j_\ell}.
\]
This counting identity follows because there are $p$ possible
positions for a marked occurrence and each word is counted $\ell$
times. The same expansion in $cx^\top R^{p-1}x=cr_{p+1}=a_p$
gives the identical sum without the factor $p/\ell$.
Every factor is nonnegative and $1\leq p/\ell\leq p$, proving
both magnitude bounds.

For the variance form diagonalize $R$, with eigenvalues $\lambda_i(R)>0$
and spectral weights $w_i$ of $b$. Put mass $cw_i \lambda_i(R)$ at $\lambda_i(R)$
and mass $1-cr_1$ at zero. These masses sum to one; in the soft
case the zero mass equals $\beta^{-1}/(\beta^{-1}+r_1)$, and it
is zero in the hard case. The resulting distribution $\nu$ has
$\E_\nu Z^j=a_j$. Thus $a_1^2\leq a_2$, and the formula for
$q_2$ from the scalar identities becomes \eqref{eq:q2variance}. The subtracted term
is at most half of $2a_2$, so second-order cancellation is uniformly
bounded, independently of dimension and condition number.
\end{proof}

\begin{theorem}[Polynomial-precision integer marginal estimation]
\label{thm:polymoment}
Use the notation of Proposition~\ref{prop:integerpositive}, at any
fixed shift. If $|\widehat r_j/r_j-1|\leq\zeta<1$ for
$j=1,\ldots,p+1$, the scalar logarithmic recurrence satisfies
\begin{equation}\label{eq:polymoment}
 \left|\frac{\widehat q_p}{q_p}-1\right|
 \leq \left(1+\frac{4\zeta}{1-\zeta}\right)^p-1.
\end{equation}
For $0<\eta\leq1$, choosing $\zeta\leq\eta/(16p)$ makes this
error at most $\eta$. The assertion holds for soft and hard grounding,
independently of dimension and condition number.
\end{theorem}
\begin{proof}
Work with formal power series through degree $p$. Set
\[
 F(z_{\rm ser})=1+\sum_{j\geq1}a_jz_{\rm ser}^j,\qquad
 D(z_{\rm ser})=\sum_{j\geq1}d_jz_{\rm ser}^j.
\]
The determinant identities in Proposition~\ref{prop:integerpositive}
give $F^{-1}=1-D$, with $a_j,d_j\geq0$, and hence
\begin{equation}\label{eq:renewalmoment}
 a_j=d_j+\sum_{i=1}^{j-1}a_i d_{j-i}.
\end{equation}
Also $q_p=p[z_{\rm ser}^p]\log F$, $q_p\geq a_p$, and Holder's inequality
for the probability measure constructed there gives
$a_{j_1}\cdots a_{j_\ell}\leq a_p$ whenever
$j_1+\cdots+j_\ell=p$ and $j_i\geq1$.

Put $\widehat a_j=\widehat r_{j+1}/(\beta^{-1}+\widehat r_1)$.
Positivity of the exact penalty yields
\[
 \frac{1-\zeta}{1+\zeta}\leq\frac{\widehat a_j}{a_j}
 \leq\frac{1+\zeta}{1-\zeta},\qquad
 |\widehat a_j-a_j|\leq\alpha a_j,\quad
 \alpha=\frac{2\zeta}{1-\zeta}.
\]
For $E(z_{\rm ser})=\sum_{j=1}^p(\widehat a_j-a_j)z_{\rm ser}^j$, define
$H=E/F=E(1-D)$. For $1\leq j\leq p$, identity
\eqref{eq:renewalmoment} implies
\[
 |[z_{\rm ser}^j]H|\leq\alpha\left(a_j+
          \sum_{i=1}^{j-1}a_i d_{j-i}\right)
 =\alpha(2a_j-d_j)\leq2\alpha a_j.
\]
The formal identity $\log(F+E)-\log F=\log(1+H)$ then gives
\begin{align*}
 |\widehat q_p-q_p|
 &\leq p\sum_{\ell=1}^p\frac{(2\alpha)^\ell}{\ell}
   \sum_{\substack{j_1+\cdots+j_\ell=p\\j_i\geq1}}
            a_{j_1}\cdots a_{j_\ell}\\
 &\leq a_p\sum_{\ell=1}^p\binom p\ell(2\alpha)^\ell
 =a_p((1+2\alpha)^p-1).
\end{align*}
Division by $q_p\geq a_p$ proves \eqref{eq:polymoment}.
All identities concern finitely many coefficients, so convergence
of a series or realizability of the perturbed moments is unnecessary.
Finally $\zeta\leq\eta/(16p)$ implies
$4p\zeta/(1-\zeta)\leq4\eta/15$.
Thus the bound is at most $e^{4\eta/15}-1\leq\eta$:
convexity bounds it by $\eta(e^{4/15}-1)$ on $[0,1]$, and
$e^{4/15}\leq(1-4/15)^{-1}=15/11<2$.
\end{proof}

\begin{proposition}[Higher-power sketch maintenance]
\label{prop:highermaintained}
Fix a matrix $\mathcal F$ and an integer $\ell_{\max}\geq1$, and store
$\mathcal Y_\ell=\mathcal F R_T^\ell$ for $1\leq\ell\leq \ell_{\max}$.
Before grounding $v$, put $R=R_T$, $x=Re_v$, $c=1/x_v$,
$a_i=R^ix$, and $s_i=x^\top a_i$. Define
\[
 b_0=x,\qquad
 b_t=a_t-c\sum_{i=0}^{t-1}s_i b_{t-1-i}
 \quad(1\leq t<\ell_{\max}).
\]
All levels update simultaneously by
\begin{equation}\label{eq:highermaintained}
 \mathcal Y_\ell'=\mathcal Y_\ell-c\sum_{i=0}^{\ell-1}
             \mathcal Y_{i+1}[:,v]b_{\ell-1-i}^\top.
\end{equation}
Every vector on the right is its pre-update value.
Apply this construction to $\mathcal F=\Pi_I$ and $\mathcal F=\Pi_C C$, with $d$ rows each.
The squared columns of $\Pi_I R_T^\ell$ and $\Pi_C C R_T^\ell$ estimate,
respectively, $r_{2\ell}(u,T)$ and $r_{2\ell-1}(u,T)$.
Thus moments through $p+1$ need $\ell_{\max}=\lceil(p+1)/2\rceil$.

Initialization uses at most $2d\ell_{\max}$ right-hand-side solves.
An update uses at most $\ell_{\max}$ selected-column solves,
$O(Nd\ell_{\max}^2)$ arithmetic, and $O(Nd\ell_{\max})$ storage; existing selected
columns can be reused. The formula specializes to the two
rank-one updates of Proposition~\ref{prop:maintainedsketch}
when $p=1$. For $p=2$, it additionally gives
\[
 \mathcal Y_2'=\mathcal Y_2-c\mathcal Y_1[:,v](Rx-c\|x\|^2x)^\top
                 -c\mathcal Y_2[:,v]x^\top.
\]
\end{proposition}
\begin{proof}
Let $\bar R_v=R-cxx^\top=R_{T+v}$. The generating-function calculation
in the incremental analysis gives $b_t=\bar R_v^tx$; alternatively it follows by comparing
coefficients in
\[
 (I-z_{\rm ser}\bar R_v)^{-1}x=
 \frac{(I-z_{\rm ser}R)^{-1}x}{1+c z_{\rm ser}x^\top(I-z_{\rm ser}R)^{-1}x}.
\]
The noncommutative telescoping identity
\[
 R^\ell-\bar R_v^\ell=c\sum_{i=0}^{\ell-1}R^ixx^\top \bar R_v^{\ell-1-i}
\]
follows by expanding the consecutive differences
$R^{i+1}\bar R_v^{\ell-1-i}-R^i\bar R_v^{\ell-i}$ and summing.
Multiplication by $\mathcal F$ proves \eqref{eq:highermaintained}, since
$\mathcal F R^ix=\mathcal Y_{i+1}[:,v]$ and $\bar R_v$ is symmetric.
Snapshotting all columns $\mathcal Y_\ell[:,v]$ before any update gives
an in-place implementation of this identity.

For every $T$, $R_TA_0R_T=R_T$. Therefore
\[
 \|R_T^\ell e_u\|^2=r_{2\ell}(u,T),\qquad
 \|CR_T^\ell e_u\|^2=r_{2\ell-1}(u,T).
\]
These identities prove the moment targets for isotropic sketches.
The chain $a_0,\ldots,a_{\ell_{\max}-1}$ uses $\ell_{\max}$ solves; its inner
products and the $b_t$ recurrence cost $O(N\ell_{\max}^2)$.
There are $\sum_{\ell=1}^{\ell_{\max}}\ell=O(\ell_{\max}^2)$ outer products per
sketch family, costing $O(Nd\ell_{\max}^2)$. The stored levels and
selected chains take $O(Nd\ell_{\max})$ space. Repeatedly solving each
of the two families of $d$ starting rows gives initialization.
For $\ell_{\max}=2$, $b_1=Rx-c\|x\|^2x$, giving the final formula.
\end{proof}



\section{Spectral Translation and Finite Removal}\label{app:spectral}

All spectral comparisons here leave a nonempty grounded matrix.
The common seed and budget fix the comparison family.
For the bicriteria statement, the certified lazy rule is described
in Appendix~\ref{app:variants}; its step factor is $1-\theta$.
The notation $X_S$ denotes the active grounded matrix $A_S^\infty$.
\begin{theorem}[Uniform inverse-power spectral rate]\label{thm:spectralrate}
Let $X\succ0$ have dimension $D$ and ordered eigenvalues
$\lambda_1\leq\cdots\leq\lambda_D$. For every $s>0$,
\begin{equation}\label{eq:spectralrate}
 \lambda_1^{-1}\leq(\tr X^{-s})^{1/s}
          \leq D^{1/s}\lambda_1^{-1}.
\end{equation}
The relative excess is at most $e^{(\log D)/s}-1$; it is at most
$\eta$ whenever $s\geq\log D/\log(1+\eta)$.
If $\lambda_1$ has multiplicity $r<D$ and
$\varrho_{\rm spec}=\lambda_1/\lambda_{r+1}<1$, then
\[
 0\leq\log\bigl(\lambda_1(\tr X^{-s})^{1/s}\bigr)
           -\frac{\log r}{s}
 \leq\frac{D-r}{rs}\varrho_{\rm spec}^s.
\]
At a fixed cardinality and matrix dimension $D$, any set whose
power-trace objective is within a multiplicative factor $\alpha_{\rm cost}\geq1$
of its minimum satisfies
\begin{equation}\label{eq:alphaspectral}
 \lambda_{\min}(X_{\rm out})
     \geq(\alpha_{\rm cost} D)^{-1/s}\lambda_k^*,
\end{equation}
where $\lambda_k^*$ is the maximum smallest eigenvalue at that
cardinality and the same seed. In particular an exact power-trace
minimizer has factor $D^{-1/s}$.
\end{theorem}
\begin{proof}
Factor the smallest eigenvalue out of the trace:
\[
 (\tr X^{-s})^{1/s}
 =\lambda_1^{-1}
    \left(\sum_{i=1}^D(\lambda_1/\lambda_i)^s\right)^{1/s}.
\]
The sum lies between 1 and $D$, giving \eqref{eq:spectralrate}
and its stated rate. With multiplicity $r$ it lies between $r$
and $r+(D-r)\varrho_{\rm spec}^s$. Taking logarithms and using
$\log(1+x)\leq x$ proves the refined inequality. If $r=D$, the
sum equals $D$ exactly.

Let $O_\lambda$ maximize the smallest eigenvalue and $O_s$ minimize
the trace, at the same budget. Then
\begin{align*}
 \lambda_{\min}(X_{\rm out})^{-s}
 &\leq\Phi_s({\rm out})\leq\alpha_{\rm cost}\Phi_s(O_s)\\
 &\leq\alpha_{\rm cost}\Phi_s(O_\lambda)
 \leq\alpha_{\rm cost} D(\lambda_k^*)^{-s}.
\end{align*}
Taking the inverse $s$th root proves \eqref{eq:alphaspectral}.
The scalar root is used to compare minimizers and spectral values;
the greedy gain remains defined using $\Phi_s=\tr X^{-s}$.
\end{proof}

\begin{theorem}[Baseline-aware translation of the gain guarantee]\label{thm:gaintranslation}
Fix $s>0$, a seed, and $k$ additional anchors. Put
$\Phi_s^0=\Phi_s(S_0)$, $\Psi_s^*=\min_{|U|=k}\Phi_s(S_0\cup U)$,
and $C_s=\Phi_s^0/\Psi_s^*$. Let $D$ be the common final dimension
and $\lambda_k^*$ the optimal smallest eigenvalue.
If a returned set satisfies a gain factor $\gamma\in(0,1]$, put
$\rho_{\rm gain}=1-\gamma$. Then
\begin{align}
 \frac{\lambda_{\min}(X_{\rm out})}{\lambda_k^*}
 &\geq\left[\rho_{\rm gain} \Phi_s^0(\lambda_k^*)^s+\gamma D\right]^{-1/s}
 \label{eq:gaintranslation}\\
 &\geq\left[D(\rho_{\rm gain} C_s+\gamma)\right]^{-1/s}.\notag
\end{align}
Thus a contrast bound $C_s\leq C$ at the chosen exponent and
\[
 s\geq\frac{\log(D(\rho_{\rm gain} C+\gamma))}{\log(1+\eta)}
\]
certify a factor $1/(1+\eta)$ on the smallest eigenvalue.
Apply this bound on the success event of any of the stated
gain guarantees, with its corresponding $\gamma$.
Without a baseline-contrast condition, any fixed gain factor
$\gamma<1$ alone admits arbitrarily small spectral ratios, even on
a connected weighted four-vertex star and at every integer exponent
$s\geq1$.
\end{theorem}
\begin{proof}
The gain guarantee is exactly the residual-objective inequality
\[
 \Phi_s({\rm out})\leq\rho_{\rm gain} \Phi_s^0+\gamma\Psi_s^*.
\]
Since $\Psi_s^*\leq D(\lambda_k^*)^{-s}$ and
$\lambda_{\min}(X_{\rm out})^{-s}\leq\Phi_s({\rm out})$,
the first spectral bound follows. Also
$\Phi_s^0(\lambda_k^*)^s=C_s\Psi_s^*(\lambda_k^*)^s\leq C_s D$,
giving the second bound and the stated sufficient exponent.
The bound holds on the corresponding gain event.

For the final assertion, take a star with center 0, leaves 1,2,3,
and respective edge weights $b/2,b,1$, where $0<b<1$. The fixed
seed is $S_0=\{0\}$, so the current matrix is
$\diag(b/2,b,1)$. Two additional anchors are allowed.
The optimal choice $\{1,2\}$ leaves the eigenvalue 1; the choice
$\{1,3\}$ leaves the eigenvalue $b$.
Their gain ratio is exactly
\[
 \frac{(2/b)^s+1}{(2/b)^s+b^{-s}}
 =\frac{2^s+b^s}{2^s+1}\geq\frac23>1-e^{-1}
 \qquad(s\geq1).
\]
Nevertheless the spectral ratio is $b$, which tends to zero.
For a larger prescribed $\gamma<1$, replace the weight $b/2$ by
$b/m_{\rm amp}$ with integer $m_{\rm amp}\geq\max\{2,\gamma/(1-\gamma)\}$.
The gain ratio is then $(m_{\rm amp}^s+b^s)/(m_{\rm amp}^s+1)\geq\gamma$ for
every integer $s\geq1$, and the spectral ratio is still $b$.
This is a counterexample to an implication from the gain factor
alone: it distinguishes that factor from the multiplicative
trace condition in \eqref{eq:alphaspectral}.
For fixed $\rho_{\rm gain}>0$ and $\lambda_k^*>\lambda_{\min}(X_{S_0})$,
the first bound in \eqref{eq:gaintranslation} tends to
$\lambda_{\min}(X_{S_0})/\lambda_k^*$ as $s\to\infty$, by
Theorem~\ref{thm:spectralrate}. Increasing the exponent makes the
surrogate sharper while the baseline term remains explicit.
\end{proof}

\begin{corollary}[A spectral bicriteria guarantee]\label{cor:spectralbudget}
Compare with the optimal set of $k$ additional anchors, with
remaining dimension $D_k$. Let the initial matrix have dimension
$D_0$, let $a>0$ lower-bound its smallest eigenvalue, and let
$B$ upper-bound $\lambda_k^*$. Run the same certified lazy rule
with step factor $1-\theta$ for $\ell\geq k$ rounds.
Whenever the budget is feasible and hard grounding leaves a vertex,
\begin{equation}\label{eq:budgetresidual}
 \Phi_s(S_\ell)
 \leq e^{-(1-\theta)\ell/k}\Phi_s^0+\Psi_{s,k}^*.
\end{equation}
For a target $\eta>0$, choose an integer
\begin{align}
 s&\geq\left\lceil\frac{\log(2D_k)}{\log(1+\eta)}\right\rceil,
 \label{eq:budgetchoices}\\
 \ell&\geq\max\left\{k,
 \left\lceil\frac{k}{1-\theta}
       \left(\log\frac{D_0}{D_k}+s\log\frac Ba\right)\right\rceil
              \right\}.\notag
\end{align}
Then, on the same probability-$1-\delta$ event,
$\lambda_{\min}(X_{S_\ell})\geq\lambda_k^*/(1+\eta)$.
\end{corollary}
\begin{proof}
The maximum-marginal argument in Theorem~\ref{thm:hybridcap} may
compare every round with a fixed size-$k$ optimum, even after
more than $k$ rounds. Its residual recurrence is
\[
 \Phi_s(S_{j+1})-\Psi_{s,k}^*
 \leq\left(1-\frac{1-\theta}{k}\right)
               (\Phi_s(S_j)-\Psi_{s,k}^*).
\]
If the parenthesized gap is negative, monotonicity also implies
this inequality. Iteration and the exponential bound prove
\eqref{eq:budgetresidual}. Now $\Phi_s^0\leq D_0a^{-s}$ and
$\Psi_{s,k}^*\leq D_k(\lambda_k^*)^{-s}$.
The second choice in \eqref{eq:budgetchoices} bounds the first
term in \eqref{eq:budgetresidual} by
$D_k B^{-s}\leq D_k(\lambda_k^*)^{-s}$. Therefore
\[
 \lambda_{\min}(X_{S_\ell})^{-s}
 \leq\Phi_s(S_\ell)\leq2D_k(\lambda_k^*)^{-s}.
\]
The first choice of exponent proves the spectral factor.
The permanent-key argument covers all $\ell$ rounds on its initial event.
\end{proof}

\begin{proposition}[Finite-removal expansion and spectral weighting]
\label{prop:zhoutaylor}
Let $X$ be a grounded Laplacian with a simple smallest eigenpair
$(\lambda_1,\phi_1)$, $\|\phi_1\|=1$, and orthonormal eigenbasis
$\phi_i$. For an active candidate $v$, let $R=X^{-1}$,
$x=Re_v$, $c=1/r_1$, $D=cxx^\top$, and $a_j=cr_{j+1}$.
For $0\leq \vartheta\leq1$ put
$Q_p(\vartheta)=\tr R^p-\tr(R-\vartheta D)^p$. Its endpoint $Q_p(1)$ is the
hard-grounding marginal. For $p=2,3,4$,
\begin{align*}
 Q_2(\vartheta)&=2\vartheta a_2-\vartheta^2a_1^2,\\
 Q_3(\vartheta)&=3\vartheta a_3-3\vartheta^2a_1a_2+\vartheta^3a_1^3,\\
 Q_4(\vartheta)&=4\vartheta a_4-\vartheta^2(4a_1a_3+2a_2^2)
                       +4\vartheta^3a_1^2a_2-\vartheta^4a_1^4.
\end{align*}
Their first-order coefficient satisfies
\[
 Q_p'(0)=\frac p{r_1(v)}\sum_i\lambda_i^{-(p+1)}\phi_i(v)^2.
\]
For the weighted Zhou score
$h_v=2\phi_1(v)\sum_{w\ne v}w_{vw}\phi_1(w)$,
\[
 h_v=2(X_{vv}-\lambda_1)\phi_1(v)^2.
\]
Whenever $X_{vv}>\lambda_1$,
\[
 Q_p'(0)=\frac{p\lambda_1^{-(p+1)}h_v}
                    {2r_1(v)(X_{vv}-\lambda_1)}
       +\frac p{r_1(v)}\sum_{i\geq2}
                          \lambda_i^{-(p+1)}\phi_i(v)^2.
\]
The tail sum is at most
$(1-\phi_1(v)^2)\lambda_2^{-(p+1)}$.
For every integer $p\geq2$,
\[
 0\leq \vartheta Q_p'(0)-Q_p(\vartheta)
 \leq\frac{p(p-1)}2 \vartheta^2 a_1^2\|R\|^{p-2}.
\]
For an infinitesimal soft penalty the coefficient is
\[
 p\sum_i\lambda_i^{-(p+1)}\phi_i(v)^2.
\]
\end{proposition}
\begin{proof}
Expand the trace of $(R-\vartheta D)^p$ as noncommutative words and use
cyclicity of the trace. Since $D$ has rank one,
$\tr(DR^{j-1})=a_j$ and
$\tr(DR^iDR^h)=a_{i+1}a_{h+1}$.
For degree four the six words with two $D$'s have traces
$4a_1a_3+2a_2^2$; the four with three $D$'s give
$4a_1^2a_2$. The degree-two and degree-three counts give the
remaining formulas. Differentiating at zero gives $pa_p$ and
spectral decomposition gives its eigenmode expression. The
$v$-coordinate of $X\phi_1=\lambda_1\phi_1$ proves the identity
for $h_v$, and substitution gives the comparison. The tail bound
uses $\sum_{i\geq2}\phi_i(v)^2=1-\phi_1(v)^2$.

Differentiating with $K_\vartheta=R-\vartheta D\succeq0$ gives
\begin{align*}
 Q_p'(\vartheta)&=p\tr(DK_\vartheta^{p-1}),\\
 -Q_p''(\vartheta)&=p\sum_{i=0}^{p-2}\tr(DK_\vartheta^iDK_\vartheta^{p-2-i}).
\end{align*}
Write $D=a_1ww^\top$ with $\|w\|=1$. Each summand equals
$a_1^2(w^\top K_\vartheta^iw)(w^\top K_\vartheta^{p-2-i}w)$, between zero
and $a_1^2\|R\|^{p-2}$ since $0\preceq K_\vartheta\preceq R$.
Integrating twice gives the remainder bound. Finally
$R-\vartheta D=(X+\beta e_ve_v^\top)^{-1}$ for
$\beta=\vartheta/((1-\vartheta)r_1)$ when $\vartheta<1$, by Sherman--Morrison.
Thus $dz/d\beta=r_1$ at zero and the soft coefficient follows.
\end{proof}

\begin{proposition}[A finite-removal stress case]\label{prop:zhoustress}
Let $0<t<1/4$. On vertices $\{z,a,b,c\}$ put edges $za,ab$
of weight $t$ and $zb,bc,cz$ of weight one. Ground the prescribed
seed $z$ and choose one further vertex. Apply the weighted version
of the score in Zhou et al.\ \cite{zhou2023}, Eq.~(4),
$h_v=2u_v\sum_{w\notin S}w_{vw}u_w$, using the exact smallest
eigenvector. The score uniquely selects $b$, and its achieved
smallest eigenvalue divided by the optimum is at most $2t$.
For every real $p\geq1$, exact minimization of $\tr X^{-p}$
over the three one-vertex choices uniquely selects the spectral
optimum $a$.
\end{proposition}
\begin{proof}
In the order $a,b,c$, the current grounded matrix is
\[
 X=\begin{pmatrix}2t&-t&0\\-t&2+t&-1\\0&-1&2\end{pmatrix}.
\]
It is positive definite by grounded connectivity. Its off-diagonal
graph is connected, so its smallest eigenvector is strictly positive:
apply Perron--Frobenius to $dI-X$ with $d>\lambda_{\max}(X)$.
The scores obey
\[
 h_a=2tu_au_b>0,\quad h_c=2u_bu_c>0,\quad h_b=h_a+h_c.
\]
Consequently $b$ is the unique maximizer, independently of how
accurately the positive eigenvector is scaled. After selecting $b$
the matrix is $\diag(2t,2)$ and its smallest eigenvalue is $2t$.
After selecting $a$, the matrix is
\[
 \begin{pmatrix}2+t&-1\\-1&2\end{pmatrix}
 =\begin{pmatrix}2&-1\\-1&2\end{pmatrix}
   +\diag(t,0)\succeq I,
\]
so its smallest eigenvalue is at least one. After selecting $c$,
the Rayleigh quotient at the coordinate $a$ is $2t<1$.
This proves uniqueness of the spectral optimum and the ratio bound.
The inverse-power trace after selecting $a$ is at most two; after
either alternative it is at least $(2t)^{-p}>2$ for every $p\geq1$.
The inverse-power choice is therefore also uniquely $a$.
\end{proof}

The opposite spectral ordering also occurs within the same family.
Consider the grounded matrix
\[
 X=\begin{pmatrix}1&-9/10&0\\-9/10&3&-1/2\\0&-1/2&1\end{pmatrix}.
\]
Its three row sums are $1/10,8/5,1/2$, so it is realized by a
connected weighted graph with one seed attached by those weights.
The follower path and positivity imply that Zhou uniquely chooses
the middle vertex. Its resulting matrix is $I_2$, giving spectral
value one and inverse trace two. Deleting the first vertex gives
eigenvalues $2\pm\sqrt5/2$, inverse trace $16/11$, and inverse
power traces $168/121,1984/1331,24352/14641$ for $p=2,3,4$,
each below two. Deleting the third leaves a smallest eigenvalue
$2-\sqrt{181}/10<7/10$, so every $p\geq2$ inverse-power trace
then exceeds two. Thus each of these three exact power-greedy
choices improves CFCC but has spectral value $2-\sqrt5/2<1$.
The two explicit examples realize distinct orderings of the two metrics.



\section{Analyzed Oracle and Selection Variants}\label{app:variants}

This appendix develops the analyzed alternatives to the implemented
CFCC hybrid: integer fresh-round selection, positive quadrature for
noninteger powers, and permanently valid lazy bounds.

\subsection{Fresh-round integer sketches}
At each current SDDM matrix $A=C^\top C$, set $P=ZC$ for a new
Gaussian $Z$, and draw an independent new Gaussian $Q$.
Squared columns of $QA^{-\ell}$ and $PA^{-\ell}$ have targets
$e_v^\top A^{-2\ell}e_v$ and $e_v^\top A^{-(2\ell-1)}e_v$.
For a fixed unit vector $w$, the Gaussian concentration bound is
\begin{equation}\label{eq:chisquare}
\Prob\{|\|Qw\|^2-1|>x\}\leq2e^{-dx^2/8},\quad 0<x<1.
\end{equation}
It is the Gaussian case proved in Theorem~\ref{thm:adaptivehybridjl}.
For approximate chains, the Frobenius error satisfies
$E_\ell\leq(E_{\ell-1}+\rho_\ell)/a$, as proved in
Proposition~\ref{prop:sketchresidual}. Expanding the perturbed
column-norm square converts this into moment error.
The SDDM primitive used below solves to relative energy precision
$\xi$ in $\widetilde O((m_E+n)\log(1/\xi))$ work, with
logarithmic failure-probability factors~\cite{spielman2014}.
Here $m_E=m$ is the input edge count, upper-bounding every round's
active off-diagonal count; diagonal entries contribute $O(n)$.
\paragraph{FreshRoundGreedy (all-node worst-case variant).}
Fix a seed $S_0$, $k$ additional vertices, integer $p\geq1$,
$0<\epsilon\leq1/2$, and $0<\delta<1$. Write $n$ for an upper
bound on the active dimension and set
\begin{equation}\label{eq:freshparams}
\begin{gathered}
 \eta=\epsilon/2,\quad\zeta=\epsilon/(32p),\quad
 J=p+1,\quad \ell_{\max}=\lceil J/2\rceil,\\
 d=\left\lceil128\zeta^{-2}\log\frac{12nJk}{\delta}\right\rceil.
\end{gathered}
\end{equation}
At each of the $k$ rounds:
\begin{enumerate}
\item Form the current SDDM matrix $A$ and its sparse factor
$A=C^\top C$ using the restricted incidence factor. Draw entirely new independent Gaussian
families $Q,Z$, with variance $1/d$ entries, and put $P=ZC$.
These families and the solver randomness are independent of all
previous rounds.
\item Apply repeated solves to these families through depth
$\ell_{\max}$, obtaining every $r_{j,v}=e_v^\top A^{-j}e_v$, $1\leq j\leq J$,
with relative tolerance $\zeta$. Specifically, squared columns of
$QA^{-\ell}$ estimate order $2\ell$, and those of $PA^{-\ell}$
estimate order $2\ell-1$. Discard orders above $J$.
\item Evaluate $\widehat\Delta_v=\widehat q_{p,v}(0)$ for every
eligible vertex by the scalar logarithmic recurrence, select a
maximizer, and ground it. Discard this round's sketches.
\end{enumerate}
Each round's conditional failure budget is $\delta/k$, including
all its candidate estimates and all its solves.

\begin{theorem}[Fresh-round all-node greedy guarantee]\label{thm:freshgreedy}
FreshRoundGreedy returns $k$ additional vertices $U_k$ with
\begin{equation}\label{eq:freshgain}
\begin{aligned}
 &\Phi_p(S_0)-\Phi_p(S_0\cup U_k)\\
 &\quad\geq(1-e^{-1}-\epsilon)
   \max_{|U|\leq k}\bigl(\Phi_p(S_0)-\Phi_p(S_0\cup U)\bigr)
\end{aligned}
\end{equation}
with probability at least $1-\delta$, for hard and soft grounding.
\end{theorem}
\begin{proof}
Condition on the complete history before round $i$. Its matrix,
factor, and eligible vertices are now fixed. The new Gaussian
families have the covariance and chi-squared laws stated above. There are at most $nJ$ moment estimates;
\eqref{eq:chisquare}, with deviation $\zeta/4$, and
\eqref{eq:freshparams} give conditional JL failure at most
\[
 2nJ\exp(-d\zeta^2/128)\leq\delta/(6k).
\]
The residual/norm budgets in Theorem~\ref{thm:freshcost} give
conditional failure at most $\delta/(3k)$ for each of the Gaussian
norm and solver events. Their intersection implies relative moment
error at most $\zeta$, by the norm-perturbation bound above. In particular the total conditional
failure is at most $\delta/k$, for every history.

Let $E_i$ be this round's simultaneous event. Taking conditional
expectations and then a union bound gives
\[
 \Prob\left(\bigcup_{i=1}^k E_i^c\right)
 \leq\sum_{i=1}^k\E\bigl[\Prob(E_i^c\mid\mathcal H_{i-1})\bigr]
 \leq\delta.
\]
Independence between the events $E_i$ themselves is unnecessary;
it is the fresh random draws conditional on the history that
justify each probability bound.

On their intersection, Theorem~\ref{thm:polymoment} gives
$(1-\eta)\Delta_v\leq\widehat\Delta_v\leq(1+\eta)\Delta_v$
for all eligible vertices in every round. If $v$ maximizes the
estimate and $u$ maximizes the true marginal, then
\[
 \Delta_v\geq\frac{\widehat\Delta_v}{1+\eta}
 \geq\frac{1-\eta}{1+\eta}\Delta_u
 \geq(1-\epsilon)\Delta_u.
\]
The normalized seeded gain is monotone and submodular by Theorem~\ref{thm:cm}.
The same gap-closing induction as in Theorem~\ref{thm:hybridcap}
therefore gives the stronger factor
$1-(1-(1-\epsilon)/k)^k\geq1-e^{-1}-\epsilon$.
For hard grounding the last allowed deletion, if any, uses the
one-dimensional marginal and leaves zero trace; no inverse of
an empty matrix is formed.
\end{proof}

\begin{theorem}[Worst-case cost with explicit polynomial dependence]
\label{thm:freshcost}
Suppose certified bounds $0<a\leq\lambda_{\min}(A)$ and
$\lambda_{\max}(A)\leq B$ hold for every round, and put $\kappa=B/a$.
Using the nearly-linear SDDM primitive stated above, FreshRoundGreedy has
\begin{equation}\label{eq:freshcost}
 T=\widetilde O\bigl(p^4 k(m_E+n)\epsilon^{-2}\bigr).
\end{equation}
Thus one explicit choice is $f(p)=p^4$; on connected graphs
$m_E+n=O(m_E)$ gives the
$\widetilde O(f(p)km_E\epsilon^{-2})$.
The suppressed factors are logarithmic in
$n,k,p,\kappa,1/\epsilon,1/\delta$, with no suppressed polynomial
factor in $p$. The count of SDDM right-hand-side solves alone is
$O(kdL)=\widetilde O(p^3k\epsilon^{-2})$.
These are input-wise worst-case bounds, independent of the observed
CELF refresh count, in the standard capped high-probability solver
model; solver failure is included in $\delta$.
\end{theorem}
\begin{proof}
Normalize every matrix by the common $B$, and also replace the
soft penalty by $\beta/B$. All original marginal values acquire
the same positive factor $B^{-p}$ on converting back, so every
selection is unchanged. The normalized spectrum is in
$[1/\kappa,1]$. Its factor satisfies
$\E\|P\|_F^2=\tr A\leq n$ and $\E\|Q\|_F^2\leq n$.
Set
\begin{equation}\label{eq:freshprecision}
 b_{\rm rnd}=\sqrt{6nk/\delta},\qquad
 \xi=\frac{\zeta}{64\ell_{\max}\kappa^{2\ell_{\max}}(b_{\rm rnd}+1)}.
\end{equation}
Markov's inequality gives simultaneous starting norms at most $b_{\rm rnd}$
except with conditional probability $\delta/(3k)$. Use relative
energy accuracy $\xi$ at every inverse-power solve, with per-solve
conditional failure at most $\delta/(6kdL)$. There are at most
$2d\ell_{\max}$ such solves per round. The norm test and a capped solve budget
may terminate with an arbitrary feasible answer on failure, so the
work bound also holds on unsuccessful executions.

For either family, let $V_0$ be its starting batch and $\widehat V_j$
its approximate chain. Write $\rho_j=\|\widehat V_j A-
\widehat V_{j-1}\|_F$ and
$E_j=\|\widehat V_j-V_0A^{-j}\|_F$. Relative energy accuracy
implies $\rho_j\leq\xi\sqrt\kappa\|\widehat V_{j-1}\|_F$.
Inductively $\|\widehat V_{j-1}\|_F\leq
\kappa^{j-1}(b_{\rm rnd}+1)$, whence for $j\leq \ell_{\max}$
\begin{gather*}
 \rho_j\leq\frac{\zeta}{64\ell_{\max}\kappa^{\ell_{\max}+1/2}}
 \leq\frac{\zeta}{64\ell_{\max}\kappa^{\ell_{\max}}},\\
 E_j\leq\kappa(E_{j-1}+\rho_j)\leq\zeta/64.
\end{gather*}
The last inequality follows by summing at most $\ell_{\max}$ terms, each
amplified by at most $\kappa^{\ell_{\max}}$, starting with $E_0=0$.
It also closes the stated induction on approximate batch norms.
Every normalized target moment is at least one. Thus these
absolute column errors, combined with the relative $\zeta/4$
JL event, meet the sufficient $\zeta/8$ norm budgets and give
relative moment error at most $\zeta$.

There are at most $2d\ell_{\max}$ solves per round, each on a matrix with
$O(m_E+n)$ nonzeros. Sparse Gaussian-factor multiplication and
moment accumulation cost $O(d(m_E+nL))$ per round. The scalar
recurrences cost $O(np^2)$ and the maximization costs $O(n)$.
The precision satisfies
\[
 \log(1/\xi)=2\ell_{\max}\log\kappa+
       \log\bigl(64\ell_{\max}(b_{\rm rnd}+1)/\zeta\bigr).
\]
A nearly-linear SDDM solve costs
$\widetilde O((m_E+n)\log(1/\xi))$, including logarithmic
failure-probability factors. Since $d=\widetilde O(p^2\epsilon^{-2})$
and $\ell_{\max}\leq p$, explicitly counting this precision cost gives
\[
 \widetilde O\bigl(k(m_E+n)p^2\ell_{\max}^2\epsilon^{-2}\bigr)
       +O(knp^2)
 \subseteq\widetilde O\bigl(p^4k(m_E+n)\epsilon^{-2}\bigr).
\]
The $p^3$ solve count and the extra factor from solve precision are
therefore distinguished. The argument applies to every input and
every adaptive matrix sequence on the simultaneous success event.
\end{proof}


\subsection{Positive scalar oracles and quadrature}
Use the active matrix and penalty notation of
Lemma~\ref{lem:resolvent}. Higher inverse powers of shifted
matrices give a positive representation for every exponent.
\begin{proposition}[Positive scalar representation for all powers]\label{prop:allpowers}
Define, for integers $m\geq1$,
\[
 q_m(t)=\frac{(-1)^{m-1}}{(m-1)!}\frac{d^{m-1}}{dt^{m-1}}q_1(t).
\]
Then $q_m(t)=\tr(A+tI)^{-m}-\tr(B+tI)^{-m}\geq0$.
For any $s>0$ and any integer $m>s$,
\begin{equation}\label{eq:allpowers}
 \Delta_s=\frac{\Gamma(m)}{\Gamma(s)\Gamma(m-s)}
       \int_0^\infty t^{m-s-1}q_m(t)dt.
\end{equation}
For integer $s=m$ one can instead use $\Delta_m=q_m(0)$.
Each $q_m(t)$ is determined by the scalar moments
$r_1(t),\ldots,r_{m+1}(t)$ from the original matrix $A$.
\end{proposition}
\begin{proof}
Repeated differentiation of the shifted inverse gives
$d^{m-1}(A+tI)^{-1}/dt^{m-1}=(-1)^{m-1}(m-1)!(A+tI)^{-m}$,
which proves the trace identity. For a soft update the ordered
eigenvalues of $A_\beta$ are at least those of $A$, by the min--max
principle. For a hard update, writing $d=\dim A$, interlacing gives
$\lambda_i(A)\leq\lambda_i(C)$ for $1\leq i\leq d-1$.
Because $x\mapsto(x+t)^{-m}$ decreases, in the soft case every matched
eigenvalue contribution is nonnegative; in the hard case this also
leaves the positive contribution of $\lambda_d(A)$. This proves
$q_m\geq0$.

The substitution $t=xu$ in the beta integral gives
\[
 \int_0^\infty t^{m-s-1}(x+t)^{-m}dt
 =x^{-s}\frac{\Gamma(m-s)\Gamma(s)}{\Gamma(m)}.
\]
Apply it to the finite eigenvalue sums and subtract. The integral
of the difference is finite: near zero its integrand is bounded
by a constant times $t^{m-s-1}$, and near infinity by a constant
times $t^{-s-1}$. This proves \eqref{eq:allpowers} as a positive
representation, including $s\geq1$.

Here is an explicit scalar evaluation rule. Put
\[
 F(t)=\begin{cases}1+\beta r_1(t),&\beta<\infty,\\r_1(t),&\beta=\infty.
 \end{cases}
 \qquad q_1(t)=-\frac{d}{dt}\log F(t).
\]
The latter identity follows immediately from \eqref{eq:qone} and
$r'_1=-r_2$. In a neighborhood of $z=0$,
\begin{align*}
 r_1(t+z)&=\sum_{j\geq0}(-1)^j r_{j+1}(t)z^j,\\
 F(t+z)&=\sum_{j\geq0}f_jz^j,\\
 \log F(t+z)&=g_0+\sum_{j\geq1}g_jz^j.
\end{align*}
The resolvent series converges for $|z|<\lambda_{\min}(A)+t$;
$F(t)>0$ also ensures a local logarithm. For soft grounding
$f_0=1+\beta r_1$ and $f_j=\beta(-1)^jr_{j+1}$ for $j\geq1$;
for hard grounding $f_j=(-1)^jr_{j+1}$ for all $j\geq0$.
Equating coefficients in $F'=F(\log F)'$ gives, successively,
\begin{equation}\label{eq:logrecurrence}
 g_k=\frac{kf_k-\sum_{j=1}^{k-1}jg_jf_{k-j}}{kf_0},\qquad
 q_m(t)=(-1)^m m g_m.
\end{equation}
Thus moments through order $m+1$ suffice. For example,
\[
 q_2(t)=\begin{cases}
 \displaystyle\frac{2\beta r_3}{1+\beta r_1}
       -\frac{\beta^2r_2^2}{(1+\beta r_1)^2},&\beta<\infty,\\[6pt]
 \displaystyle\frac{2r_3}{r_1}-\frac{r_2^2}{r_1^2},&\beta=\infty.
 \end{cases}
\]
Its nonnegativity, and that of every higher-order expression, follows
from the trace/interlacing argument above.
\end{proof}


For the quadrature proof, let $N_A(x)$ and $N_B(x)$ count
active eigenvalues at most $x$, and put $\xi(x)=N_A(x)-N_B(x)$.
The rank-one and principal-submatrix interlacing inequalities give
$0\leq\xi\leq1$. If both spectra lie in $[a,B_*]$, integration
of each scalar eigenvalue contribution gives
$\Delta_f=\int_a^{B_*}(-f')\xi$ for soft grounding and
$\Delta_f=f(B_*)+\int_a^{B_*}(-f')\xi$ for hard grounding.
The next statement calls the spectral upper bound $B$.
\begin{theorem}[Fixed positive surrogate with logarithmic shift count]\label{thm:fixedquad}
Suppose every nonempty admissible grounded matrix has spectrum in
$[a,B]$, $a>0$, and put $\kappa=B/a$. For each $s>0$ and
$0<\eta<1$, there is a fixed positive rule
\[
 \widetilde f_s(x)=\sum_{j=1}^q w_j(x+t_j)^{-m},\quad
 t_j,w_j>0,\quad m=\lfloor s\rfloor+1,
\]
such that its trace objective is nonincreasing and supermodular at
every rule order and, for every soft or hard marginal,
\begin{equation}\label{eq:fixedrelative}
 (1-\eta)\Delta_{x^{-s}}\leq\Delta_{\widetilde f_s}\leq\Delta_{x^{-s}}.
\end{equation}
One explicit choice is
\[
 p=m+3+\left\lceil\log_4(2/\eta)\right\rceil,\qquad
 q\leq(\lceil\log_2\kappa\rceil+2)p.
\]
Thus $q=O_s((1+\log\kappa)\log(1/\eta))$ for a fixed exponent.
For $0<s<1$ the rule uses only first inverse powers ($m=1$);
for $s=1$ one can use the exact single shift $t=0$ instead.
\end{theorem}
\begin{proof}
Use the positive integral in Proposition~\ref{prop:allpowers}.
Partition its shift domain into $[0,a]$, geometric intervals
$[a,2a],[2a,4a],\ldots$ truncated at $B$, and $[B,\infty)$.
On $[0,a]$ use $t=az$ and $p$-node Gauss--Jacobi for the weight
$z^{m-s-1}$ on $[0,1]$. On every interior interval use $p$-node
Gauss--Legendre, including the positive factor $t^{m-s-1}$ in
the sampled integrand. On the tail use $z=B/t$ and Gauss--Jacobi
with weight $z^{s-1}$. The tail integrand becomes
$B^{-s}(1+xz/B)^{-m}$, a positive multiple of
$(x+B/z)^{-m}$. All nodes and coefficients are positive and
independent of $S$ or the candidate. Write the resulting scalar
approximation as $F_p(x)$.

Here is a uniform error bound for both $f_s$ and $g_s=-f_s'$.
Put $\alpha=s+1-m\in[0,1)$ and $D=2p-1$.
For an interior interval with center $c$ write $y=(t-c)/c$;
its doubling ratio ensures $|y|\leq1/3$. After extracting a
positive factor, the integrands for the function and its negative
derivative have the form
$(1+y)^{-\alpha}(1+\theta y)^{-r}$ with $0\leq\theta\leq1$
and $r=m,m+1$. On $|y|=2/3$ their modulus is at most
$3^{\alpha+r}$, while their minimum on the real interval is at
least $(4/3)^{-\alpha-r}$.
Cauchy's coefficient bound shows that the Taylor tail beyond degree
$D$ is at most $2\,3^{\alpha+r}2^{-(D+1)}$.
The quadrature integrates the Taylor polynomial exactly; positivity
and exactness on constants bound its error plus the integral's error
by twice this remainder times the interval mass. Dividing by the
true integral gives relative error at most
$4\,4^{\alpha+r}4^{-p}\leq4^{m+3-p}$.

For the lower endpoint the weight absorbs the singularity, leaving
$(x+az)^{-r}$. Expanding about $z=1/2$ gives a normalized variable
of radius $a/(2x+a)\leq1/3$ and the same, or smaller, bound.
For the upper endpoint the derivative integrand, under its fixed
weight, is proportional to $z(1+xz/B)^{-m-1}$.
Approximate its denominator factor to degree $D-1$ and multiply
by $z$. The error argument now uses exactness on $z$, with radius
$x/(2B+x)\leq1/3$. The bound is
$8\,4^{m+1}4^{-p}\leq4^{m+3-p}$.
The function tail needs degree $D$ and has a smaller bound.
These estimates hold for every $x\in[a,B]$. Summing positive panel
integrals gives, with $\epsilon=4^{m+3-p}\leq\eta/2$,
\[
 |F_p/f_s-1|\leq\epsilon,\qquad
 |(-F_p')/g_s-1|\leq\epsilon.
\]
Define $\widetilde f_s=F_p/(1+\eta/2)$. Both it and its negative
derivative lie between $1-\eta$ and 1 times their true counterparts.

For soft grounding integrate the derivative bounds against the
nonnegative spectral count $\xi$. For hard grounding use
$\Delta_f=f(B)+\int_a^B(-f')\xi$, so the endpoint value is controlled
as well. This proves \eqref{eq:fixedrelative} without taking a
relative error of two nearly equal traces. In particular it treats
the hard ratio $r_2/r_1$ rather than a log-resolvent substitute.
Every shifted inverse power is CM, so Theorem~\ref{thm:cm} proves
supermodularity of every finite positive rule, including orders
below the accuracy threshold. There are at most
$\lceil\log_2\kappa\rceil+2$ panels, proving the count.
\end{proof}


\subsection{All-node shifts and certified lazy evaluation}
For a fixed shifted system $R_t=(A+tI)^{-1}$, take $P=ZC$,
$Q$ independent, and reuse these right-hand sides across shifts.
The squared columns of $QR_t$ estimate the second resolvent
moment; those of $(P+\sqrt tQ)R_t$ estimate the first, since
$R_t(A+tI)R_t=R_t$. This finite-shift construction avoids an
additional integral for the first moment. For a common right-hand
side, unpreconditioned shift-invariant Krylov spaces provide another
implementation of the same shifted solves. Its work retains the
Krylov depth explicitly.

For deterministic refinement, a single Lanczos run yields Gauss
lower and left-Radau upper bilinear bounds for the CM function
$-f'$. Integrating these bounds along the grounding penalty gives
marginal intervals. Their one-sided signs follow from the even
and odd Hermite remainders, respectively; finite hard-tail bounds
control the infinite penalty endpoint.

A certified lazy rule starts with valid marginal upper keys.
At the current heap top $v$, refine an interval
$L_v\leq\Delta_v\leq H_v$ until
$L_v\geq(1-\theta)H_v$. Replace its key by the minimum of the
old and fresh upper bounds; accept if
$L_v\geq(1-\theta)U_{\rm next}$, otherwise reinsert it.
Supermodularity preserves all old upper bounds. A rejected fresh
candidate falls strictly below another key; when it reaches the
top again in the same round, its relative-width condition makes
it pass. Thus each candidate is refreshed at most once per round.
Every accepted marginal is at least $(1-\theta)$ of the current
maximum. The gap induction in Theorem~\ref{thm:hybridcap} gives
$1-(1-(1-\theta)/k)^k\geq1-e^{-1}-\theta$ for seeded gain.
The same induction with a fixed size-$k$ comparison set proves
the augmented-budget statement in Corollary~\ref{cor:spectralbudget}.

For general heat/CM refinement, the required Lanczos depth also
depends on the chosen function and, for heat, its horizon.
The integer fresh-round cost and the CFCC hybrid accounting use
their respective stated oracle models.


\section{Implementation Details}\label{app:implementation}
The engine stores both solution sketches transposed, as $N\times d$ arrays with
unscaled $\pm1$ entries generated by a seeded hash of the row index
and the edge (for $\Pi_C$) or vertex (for $\Pi_I$) index, and divides squared norms by $d$; this is equivalent to
the variance-$1/d$ convention of \eqref{eq:banks}. A candidate is
marked \emph{refreshed} once its moments have been computed by numerical
refresh solves; the analysis assigns those solves the explicit error
budgets of Appendix~\ref{app:numerical}. The round heap orders equal
keys by decreasing vertex index, and only a strictly larger refreshed
value replaces the incumbent. A reused refresh column keeps its
refresh tolerance $\varepsilon_{\rm ref}$. The approximate-Cholesky factor is applied
to blocks of at least four right-hand sides, and Jacobi-preconditioned
CG is used for narrower blocks; on roadNet-PA the factor is applied to
every block, including single right-hand sides. The sketch-only
ablation keys every candidate from the solution sketches in every round, and the
scalar-cache ablation keeps no solution sketches.

\clearpage
\section{Additional Experimental Details}\label{app:extraresults}

The shared seed stage of the main comparison is listed in
Table~\ref{tab:seeds}. These costs are common to all compared
selection methods. Adding the same seed cost to both methods
preserves their ordering while changing an end-to-end speed ratio.

\begin{table}[H]
\centering\footnotesize
\caption{Shared JL seed-stage times, dimension 128.}
\label{tab:seeds}
\begin{tabular}{@{}lr@{}}\toprule
Graph & Seed time (seconds)\\\midrule
DBLP & 3.747\\
Amazon & 5.400\\
YouTube & 12.710\\
web-Google & 12.972\\
Skitter & 21.157\\
roadNet-PA & 25.343\\
LiveJournal & 62.708\\
\bottomrule\end{tabular}
\end{table}

